\section{Prompts, requests and provenance}\label{app:provenance}

\paragraph{Scenes and prompts.}
The 14 scene descriptions (Table~\ref{tab:scenes}) are a fixed authored panel, not a sample of
possible content. Each prompt is the clause, a space, the scene description and ``No text or
frame.'' The named clauses read ``Render as an oil painting in the style of'' followed by
``Claude Monet.'', ``Alfred Sisley.'', ``Camille Pissarro.'' or ``Paul Cezanne.''; the
artist-free prompt has no clause.

% Generated by paper/tmlr/build_assets.py; do not edit by hand.
% input data/manifests/painter_specificity_v2/psv2-20260911/requests.jsonl sha256 0aa60f491be08dbd48f1f018e225b22bcae2aacd724abf3371a6312d59acc679
\begin{table}[h!]
\caption{The 14 scene descriptions, with their fixed content class. Each prompt is the clause, the description and ``No text or frame.''}
\label{tab:scenes}
\vspace{2pt}
\centering\small
\begin{tabular}{@{}rlp{0.78\linewidth}@{}}
\toprule
\# & Class & Scene description\\
\midrule
1 & water & A broad river bends past a gravel bank. Three tall trees stand on the far bank, with a low hill behind them under a pale overcast sky.\\
2 & water & A small pond lies between reeds and a grassy slope. A wooden landing extends from the left bank, with scattered clouds reflected on the water.\\
3 & water & A narrow canal passes beneath a low stone bridge. Garden walls line one side and leafy trees line the other in soft afternoon light.\\
4 & water & A quiet inlet meets a flat sandy shore. Two small boats rest near the water and distant headlands sit beneath a light grey sky.\\
5 & built & A village street curves past modest stone houses. A low wall and a single tree occupy the foreground, with hills visible beyond the roofs.\\
6 & built & A country house stands beside an orchard. A gravel path crosses the foreground and a small shed stands on the right under a clear morning sky.\\
7 & built & A square stone tower rises behind a cluster of low tiled roofs. An open field fills the foreground, with a few shrubs along a narrow footpath.\\
8 & built & A courtyard is enclosed by two simple farm buildings. A gate opens toward distant fields and a patch of sunlight falls across the bare ground.\\
9 & land & An open meadow slopes gently toward distant wooded hills. A narrow path crosses the grass and a few irregular clouds fill the upper sky.\\
10 & land & A rocky hillside carries scattered pine trees. A broad valley extends behind the foreground rocks, with distant mountains in hazy daylight.\\
11 & land & Rows of fruit trees stretch across a gently rising field. Tall grass borders the nearest row and a dark wood lies at the horizon.\\
12 & mixed & A river passes a small village at the foot of a hill. A stone wall crosses the foreground, with open grass and two trees beside the water.\\
13 & mixed & A farmhouse stands above a shallow stream. A footbridge connects two grassy banks and an orchard spreads behind the house in diffuse daylight.\\
14 & mixed & A lakeside path passes a low stone building and a cluster of trees. Water occupies the left half and a gentle ridge crosses the background.\\
\bottomrule
\end{tabular}
\end{table}


\paragraph{Collection history.} A first collection attempt was stopped after a technical probe
showed that an earlier gateway route accepted an invented model identifier, so that model selection
could not be verified; none of its images was measured or used. The protocol was then revised to use
explicit routes, and, before any image of the new collection existed, the scene panel was reduced
from 16 to the 14 descriptions in Table~\ref{tab:scenes} (two scenes of the original 16 were
omitted) to fit the budget at medium quality. Nothing was selected on the basis of outcomes.

\paragraph{Request configurations.}
All requests used the OpenRouter image API with a single fixed provider per configuration
(OpenAI for the four GPT Image configurations, Google AI Studio for Nano Banana 2 and Black
Forest Labs for FLUX.2 Max), fallbacks disabled, one image per request, square aspect ratio,
and no image input or seed. The OpenAI configurations requested medium quality and an opaque
background; Nano Banana 2 requested 1K resolution; FLUX.2 Max requested PNG output. The exact
payloads and a randomized request order were fixed and hashed before collection, which ran on
10 September 2026 between 15:34 and 18:21 UTC; all 1,008 requests succeeded and all images are
$1024\times1024$ pixels with distinct hashes. The responses do not echo model, quality or size
fields, so the labels identify requested configurations. The comparison concerns these
configurations, not each service's best attainable quality or equal cost per image.

\paragraph{Model selection.} The stopped first attempt had shown that an earlier route returned an
image for an invented model identifier. The explicit OpenAI routes used here were not tested with
an invalid identifier. Two observations show that the six requested configurations produced
distinguishable outputs, although they cannot establish distinct checkpoints. First, the cost per medium-quality image differed among the OpenAI
configurations in pre-collection probes (GPT Image 1 \$0.042, GPT Image 2 \$0.053, Flare and
Sunburst \$0.013 each). Second, for every pair of configurations, including Flare and Sunburst (95.2,
91.7, 100.0\% in the 31 features, CLIP and CSD), the images of the two configurations lie farther
apart than each configuration's own two repeats in most scene--clause cells, at least 91.7\% of them
for every pair and representation.

\paragraph{Reference collections.}
Eligible works were recorded in Wikidata as paintings with a single creator, the painter, and
classified as outdoor places by a lexicon applied to titles and metadata that was fixed before
any feature was measured; each needed a Wikimedia Commons reproduction at least 1,024 pixels on
its short side. Works were assigned to the development and reference panels by a fixed hash
rule. By a title-based content lexicon, the reference collections comprise (water, built place,
route, open or wooded land) 191, 67, 8 and 31 Monet works; 52, 29, 9 and 16 Sisley works; 28,
77, 12 and 24 Pissarro works; and 31, 28, 12 and 34 C\'ezanne works. Source records come from Wikimedia
Commons \citep{wikimediacommons}; the panels were assembled before the experiment.

\paragraph{Repeat dependence.}
The cross-repeat estimators assume that the two repeats of a cell have independent errors, and
two repeats cannot test this. If both repeats of every cell shared an error component with a
common fraction $\rho$ of the noise power, the cross-repeat products would include that shared
component. The squared sizes $N$ and $B$ would then be biased upward, and so would each
configuration's error $D$, by an amount proportional to its noise power. Recomputing all 15
pairwise differences under this assumption, three point orderings reverse for $\rho<1$: Nano
Banana 2 and FLUX.2 Max at $\rho=0.251$, GPT Image 1 and GPT Image 2 at $\rho=0.246$, and GPT
Image 1 and FLUX.2 Max at $\rho=0.688$. The actual dependence is not identified, so we treat the
pairwise orderings as conditional on independent repeats. The same model moves the comparisons
with 1: GPT Image 1's error would fall below 1 at $\rho=0.22$, whereas Flare's and Sunburst's stay
above 1 up to $\rho=0.64$ and 0.52. Nano Banana 2's and FLUX.2 Max's adjusted errors become negative,
which an expected squared distance cannot be, at $\rho=$ 0.30--0.32, so the point estimates imply a
dependence below about 0.3 under this model. In the embeddings (Appendix
Table~\ref{tab:embedding-calibration}), the CLIP errors above 1 would fall to 1 at $\rho=$
0.03--0.14, whereas the CSD errors, all below 1, would only fall further. Appendix
Tables~\ref{tab:robustness} and~\ref{tab:embedding-calibration} report each configuration's repeat
noise.

The request records allow one descriptive check. Scenes were requested in contiguous randomized
blocks, and the two repeats of a cell started a median of 3.1 minutes apart (at most 20 minutes).
A common linear drift in time, fitted to the 84 repeat differences of a configuration and
evaluated on held-out scenes, predicts those differences worse than no drift in every
configuration: its held-out squared prediction error exceeds that of predicting zero by 0.7--3.4\%.
This
argues against slow drift within the collection but cannot rule out nonlinear or shared-state
dependence.

\FloatBarrier
\section{Feature definitions}\label{app:features}

\begin{table}[ht]
\caption{The 31 image features. Counts are scalar coordinates. Texture features describe the
digital image; they do not measure physical paint relief.}
\label{tab:features}
\vspace{2pt}
\centering\small
\begin{tabular}{@{}p{0.14\linewidth}p{0.42\linewidth}p{0.36\linewidth}@{}}
\toprule
Family & Measurements & Interpretation\\
\midrule
Color (11) & Lightness and chroma median and IQR; chromatic fraction; hue concentration and
entropy; color differences at three lags and their slope & Brightness, color strength and
diversity, and how quickly colors change across the image\\[3pt]
Spatial (8) & Spectral slope, residual and anisotropy; orientation entropy;
horizontal--vertical balance; gradient median and IQR; quadrant divergence & Coarse versus fine
organization, dominant directions, edge strength and regional variation in orientation\\[3pt]
Texture (12) & Four wavelet energies, their slope and curvature; three local-binary-pattern
entropies; three local coefficients of variation & Detail across scales, diversity of local
patterns and local luminance contrast\\
\bottomrule
\end{tabular}
\end{table}

Images are orientation-corrected and converted to sRGB using valid embedded color profiles;
unprofiled files are treated as sRGB. They are resized with Lanczos filtering to a 512-pixel
short side without upsampling, preserving aspect ratio. No painting-region mask is applied in
the primary analysis. The features are correlated with each other; Appendix
Table~\ref{tab:sensitivity} and Table~\ref{tab:families} report alternative weightings.

\paragraph{Color.} Features use CIELAB under D65 and the $2^\circ$ observer. Lightness and
chroma each contribute a median and an IQR. The chromatic fraction is the share of pixels with
$C^*\ge5$. Hue concentration is the length of the chroma-weighted circular mean hue; hue
entropy uses 24 equal bins, normalized by the log of the bin count, with chroma weights; both
are zero when fewer than 1\% of pixels are chromatic. At lags of 1\%, 4\% and 16\% of the short
side, the median CIEDE2000 difference \citep{sharma2005ciede} is taken over pooled horizontal
and vertical pixel pairs, and a log--log slope summarizes the three.

\paragraph{Spatial.} Features use linear-sRGB luminance. After a Tukey window (parameter 0.1),
32 logarithmic radial bins cover 4--128 cycles per short side; a Theil--Sen line gives the
spectral slope and its root-mean-square residual, and power-weighted axial concentration gives
anisotropy. Scharr derivatives give the gradient median and IQR, an 18-bin magnitude-weighted
orientation entropy and the horizontal--vertical balance. Quadrant divergence is the mean
Jensen--Shannon divergence of the four quadrants' orientation histograms from the global one.

\paragraph{Texture.} A four-level undecimated \texttt{db2} wavelet transform
\citep{mallat1989wavelet} gives four log mean-squared detail energies, their linear slope and
mean second difference. Uniform local-binary-pattern entropies \citep{ojala2002lbp} use
$(P,R)=(8,1),(16,2),(32,4)$ on quantized luminance. Three median local coefficients of
variation use windows of 1\%, 4\% and 16\% of the short side (rounded up to odd widths) with a
$10^{-6}$ floor on the mean.

\paragraph{Scaling.} Each coordinate is centered by the median and divided by the IQR of the
221-work development panel (101 Monet, 36 Sisley, 48 Pissarro and 36 C\'ezanne), with each
painter given equal total weight.

\FloatBarrier
\section{Estimators}\label{app:estimators}

\begin{table}[ht]
\caption{Estimands, their meaning and their reference values. All squared sizes use
cross-repeat products.}
\label{tab:estimands}
\vspace{2pt}
\centering\small
\begin{tabular}{@{}lp{0.52\linewidth}p{0.3\linewidth}@{}}
\toprule
Symbol & Meaning & Reference values\\
\midrule
$N$, $B$ & Squared size of the shared change beyond the generic clause, and of the
between-name differences & ---\\
$N/(N+B)$ & Shared fraction & 25\% exchangeable null; $N^*/(N^*+H)$ faithful imitator;
$N/(N+H)$ exact differences\\
$H$ & Spread of the four reference painter means, $\sum_a\|r_a\|^2$ & ---\\
$B/H$ & Size of the between-name differences relative to the reference differences & 1 for a
faithful imitator\\
$\cos(c,t)$, $\lambda$ & Direction of the shared change toward the reference centroid, and the
fraction of that distance covered & 1 and 1 for a faithful imitator\\
$\beta$ & Aligned amplitude of the between-name differences & 1 faithful; 0 no painter
distinctions\\
$Q$ & Squared size of the centered generated differences relative to $H$,
$B/H+V_{\mathrm{scene}}$ & 1 faithful; 0 no distinctions\\
$\beta/\sqrt{Q}$ & Alignment ratio: alignment of the scene-wise differences with the pooled
reference pattern & 1 faithful\\
$D$, $D_{\mathrm{agg}}$ & Error of the scene-wise and of the scene-averaged differences & 0
faithful; 1 no distinctions\\
$V_{\mathrm{scene}}$ & Variation of the differences across scenes, $D-D_{\mathrm{agg}}$ & 0 if
constant across scenes\\
$D_{\mathrm{held}}$ & Error after rescaling the differences by a held-out scalar & ---\\
\bottomrule
\end{tabular}
\end{table}

\paragraph{Noise correction.} Write a repeat as $x_k=\theta+\eta_k$ with fixed mean $\theta$
and errors of mean zero that are independent between $k=1,2$. Then
$\mathbb E\ip{x_1}{x_2}=\|\theta\|^2$, whereas $\mathbb E\|x_1\|^2=\|\theta\|^2+\mathbb
E\|\eta_1\|^2$. Products of two different vectors are symmetrized,
$\frac12(\ip{u_1}{v_2}+\ip{u_2}{v_1})$. Centering across the four painters couples the painters
within a repeat but does not introduce dependence between repeats.

\paragraph{Exchangeable null.} If the four named shifts $h_a$ are independent with mean zero and
equal expected squared size $\sigma^2$, then $\mathbb E[N]=\mathbb E\,4\|\bar h\|^2=\sigma^2$ and
$\mathbb E[N]+\mathbb E[B]=\mathbb E\sum_a\|h_a\|^2=4\sigma^2$, so the ratio of expectations
$\mathbb E[N]/(\mathbb E[N]+\mathbb E[B])$ is $1/4$.

\paragraph{Baselines and the cross term.} Averaging scenes within repeat $k$, let $g_k$ be the
generic-minus-free shift, $c_k$ the shared change beyond the generic clause, and $g_k+c_k$ the
shared change against the artist-free baseline. Then
\begin{equation}
\underbrace{4\ip{g_1+c_1}{g_2+c_2}}_{N_{\mathrm{free}}}=\underbrace{4\ip{g_1}{g_2}}_{G}
+\underbrace{4\ip{c_1}{c_2}}_{N}
+\underbrace{4\bigl(\ip{g_1}{c_2}+\ip{c_1}{g_2}\bigr)}_{I}.
\end{equation}
The share of $N$ along the generic shift is $4\,\ip{c}{g}^2/(\|g\|^2N)$ with symmetrized
products. The within-scene variant computes the shared and between-name components in each
scene before averaging them.

\paragraph{Centroid proximity.} Write the named mean for painter $a$ as $\bar z_a=\bar
z_{\mathrm{g}}+c+e_a$ and let $t=\bar\mu-\bar z_{\mathrm{g}}$. Because $\sum_ae_a=\sum_ar_a=0$,
the mean reduction in squared distance to the painters' reference means is
\begin{equation}
\frac14\sum_a\bigl(\|\mu_a-\bar z_{\mathrm{g}}\|^2-\|\mu_a-\bar z_a\|^2\bigr)
=\underbrace{2\ip{c}{t}-\|c\|^2}_{\text{shared}}
+\underbrace{\tfrac12\textstyle\sum_a\ip{e_a}{r_a}-\tfrac14\sum_a\|e_a\|^2}_{\text{between-name}},
\label{eq:centroid}
\end{equation}
which holds exactly with symmetrized cross-repeat products.

\paragraph{Split of $D$.} With repeat-specific slopes $s_{sk}=\sum_a\ip{d_{sak}}{r_a}/H$ and
residuals $o_{sak}=d_{sak}-s_{sk}r_a$ orthogonal to the reference pattern, each scene's error splits
exactly as $(s_{s1}-1)(s_{s2}-1)+H^{-1}\sum_a\ip{o_{sa1}}{o_{sa2}}$: the error of the amplitude along
the pattern and the off-pattern part. The protocol prescribed this split to distinguish weak or
strong aligned responses from differences in other directions.

\paragraph{Error identities.} With $\theta_{sa}=\mathbb E[d_{sak}]$ under the repeat model,
$\mathbb E[\widehat D]=(SH)^{-1}\sum_{s,a}\|\theta_{sa}-r_a\|^2$. Writing
$\bar\theta_a=S^{-1}\sum_s\theta_{sa}$,
\begin{equation}
D=\underbrace{\frac1H\sum_a\|\bar\theta_a-r_a\|^2}_{D_{\mathrm{agg}}}
+\underbrace{\frac{1}{SH}\sum_{s,a}\|\theta_{sa}-\bar\theta_a\|^2}_{V_{\mathrm{scene}}}.
\end{equation}
Because $\beta$ is the same for the scene-averaged differences, $Q=B/H+V_{\mathrm{scene}}$.
Rescaling every centered generated contrast by $\kappa$ gives $D(\kappa)=1-2\kappa\beta+\kappa^2Q$,
minimized at $\kappa=\beta/Q$ with value $1-\beta^2/Q$; we fit $\kappa=\max(0,\beta/Q)$ on 13
scenes and score the held-out scene.

\paragraph{Finite-sample bias of $H$.} With $n_a$ reference works and within-painter covariance
$\Sigma_a$, $\mathbb E\,\widehat H=H+\frac34\sum_a\operatorname{tr}(\Sigma_a)/n_a$; subtracting
the estimated bias lowers $H$ by 6.7\%. Because $\beta$ and $Q$ are divided by $H$ and
$D-1=Q-2\beta$, all three scale by the same factor. Both benchmark fractions rise: $N/(N+H)$
because $H$ falls, and $N^*/(N^*+H)$ because the noise of the centroid $\bar\mu$ adds only a third
of the correction to $N^*$, which is much larger than $H$.

\paragraph{Intervals.} For a pair of configurations, one error difference per scene gives a
paired-scene standard error; the simultaneous half-width is $t_{13,\,1-0.05/42}$ times that
standard error, and the six aligned amplitudes use the same adjustment. Because the scenes were
authored rather than sampled, these intervals describe variation across these 14 scenes and do
not generalize to a population of scenes; for the mean over these fixed scenes, the between-scene
standard error is conservative because it also includes the variance of the scene effects. In
5,000 simulated experiments per scenario with the same
design, independent Gaussian, heavy-tailed heteroskedastic and scene-interaction errors gave
simultaneous coverage of 96.0\%, 97.4\% and 99.7\%, whereas a configuration-specific state
shared by all scenes and both repeats reduced it to 0.04\%. In that scenario, which would arise if a
closed service drifted or cached state during the collection so that both repeats of every cell
moved together, the cross-repeat products no longer remove the noise, and every interval is shifted.
Two repeats cannot rule it out, which is why repeat dependence is listed as a limitation.
Reference resampling draws 1,000 bootstrap samples of the reference works within painter, holding
the generated vectors fixed. In the first collection's analyses, scene-bootstrap draws in which a
denominator is not positive are dropped and counted; the second collection's prespecified test H1
keeps them (Appendix~\ref{app:second}).

\FloatBarrier
\section{Additional results for the 31 features}\label{app:hand}

% Generated by paper/tmlr/build_assets.py; do not edit by hand.
% input reports/painter_specificity_review_v3/analysis.json sha256 79fa5ca28900e852c072e2a52d35f3cb400ee053935d857729b4ad4aa3bc542c
% input reports/painter_specificity_review_v1/analysis.json sha256 2fabb78e15dfcc7b027a73170f6ad0c87af0196c754b4e6f45275269c2c961e2
% input reports/painter_tmlr_diagnostics_v1/analysis.json sha256 76792fa35f90e3d8fec1a75b0eec27d69cac14ed875fce8b64096b8e2525a4d4
\begin{table}[t]
\caption{Complete decomposition in the 31 features, in units of $H=5.915$. $G$: artist-free-to-generic shift. $N$: shared change beyond the generic clause. $I$: cross term. $N_{\mathrm{free}}=G+N+I$: shared change against the artist-free baseline. $B$: between-name differences (identical for both baselines). Within scene: $N/(N+B)$ with both components computed in each scene before averaging.}
\label{tab:decomposition-full}
\vspace{2pt}
\centering\small
\begin{tabular}{@{}lrrrrrrrr@{}}
\toprule
 & \multicolumn{5}{c}{Squared change $/H$} & \multicolumn{3}{c}{Shared fraction (\%)}\\
\cmidrule(lr){2-6}\cmidrule(l){7-9}
Configuration & $G$ & $N$ & $I$ & $N_{\mathrm{free}}$ & $B$ & vs.\ free & vs.\ generic & within scene\\
\midrule
GPT Image 1 & 5.04 & 5.26 & $-$0.29 & 10.01 & 2.12 & 82.5 & 71.3 & 71.0\\
GPT Image 2 & 6.28 & 4.98 & 1.06 & 12.32 & 2.04 & 85.8 & 70.9 & 76.2\\
Flare & 4.01 & 4.99 & 2.29 & 11.29 & 2.03 & 84.8 & 71.1 & 76.9\\
Sunburst & 3.08 & 3.97 & 3.28 & 10.32 & 1.98 & 83.9 & 66.7 & 72.2\\
Nano Banana 2 & 1.22 & 1.36 & 1.91 & 4.49 & 0.61 & 88.1 & 69.1 & 52.8\\
FLUX.2 Max & 3.78 & 5.34 & 6.53 & 15.65 & 0.70 & 95.7 & 88.4 & 93.6\\
\bottomrule
\end{tabular}
\end{table}

% Generated by paper/tmlr/build_assets.py; do not edit by hand.
% input reports/painter_tmlr_diagnostics_v1/analysis.json sha256 76792fa35f90e3d8fec1a75b0eec27d69cac14ed875fce8b64096b8e2525a4d4
% input reports/painter_tmlr_diagnostics_v2/analysis.json sha256 bed891f767a98f4b281b6783b1021956f4d2c3ff80becbf6b2693ee45987ef1a
% input reports/painter_cross_cohort_v1/analysis.json sha256 9488dfafa5779b8b065cc756804e67ad64c2841a5c62c2f9b0105f8fef13973e
\begin{table}[t]
\caption{Shared fraction $N/(N+B)$ (\%) of what the names add beyond the generic clause, by feature family and under two alternative weightings of all 31 features; deleting any one feature keeps it within 64.8--89.6\%. The interval row gives 95\% intervals for the mean over the six configurations from 5,000 scene resamples; the faithful and exact-differences rows average the six configurations' benchmark values in each family. The SD-Turbo row uses its own baseline, which already requests an oil painting.}
\label{tab:families}
\vspace{2pt}
\centering\small\setlength{\tabcolsep}{4pt}
\begin{tabular}{@{}lcccccc@{}}
\toprule
 & \multicolumn{4}{c}{Feature family} & \multicolumn{2}{c}{Weighting}\\
\cmidrule(lr){2-5}\cmidrule(l){6-7}
Configuration & all 31 & color & spatial & texture & equal family & covariance\\
\midrule
GPT Image 1 & 71.3 & 83.8 & 77.7 & 48.6 & 72.4 & 72.2\\
GPT Image 2 & 70.9 & 74.3 & 59.6 & 73.8 & 70.0 & 71.2\\
Flare & 71.1 & 81.5 & 65.5 & 60.1 & 70.8 & 76.4\\
Sunburst & 66.7 & 66.7 & 61.5 & 70.3 & 66.1 & 71.3\\
Nano Banana 2 & 69.1 & 56.8 & 69.8 & 79.8 & 68.9 & 65.8\\
FLUX.2 Max & 88.4 & 82.6 & 95.1 & 76.5 & 89.8 & 86.9\\
\midrule
Mean of the six & 72.9 & 74.3 & 71.5 & 68.2 & &\\
\quad scene 95\% interval & [68.8, 75.6] & [70.9, 77.2] & [57.3, 80.1] & [61.9, 70.7] & &\\
\quad faithful, mean & 89.8 & 87.2 & 85.5 & 89.4 & &\\
\quad exact differences, mean & 78.7 & 77.8 & 79.1 & 74.7 & &\\
SD-Turbo (separate) & 64.2 & 57.5 & 84.9 & 36.6 & &\\
\bottomrule
\end{tabular}
\end{table}

% Generated by paper/tmlr/build_assets.py; do not edit by hand.
% input reports/painter_tmlr_diagnostics_v1/analysis.json sha256 76792fa35f90e3d8fec1a75b0eec27d69cac14ed875fce8b64096b8e2525a4d4
\begin{table}[t]
\caption{Direction of the shared change $c$ in the 31 features, with cosines corrected by cross-repeat products. $\cos(c,t)$: cosine with the vector $t$ from the generic mean to the reference centroid. Covered: $\lambda=\langle c,t\rangle/\|t\|^2$, the fraction of $t$ covered along its direction. $\cos(c,g)$: cosine with the artist-free-to-generic shift $g$. Along generic: the share of $N$ along $g$, $\cos^2(c,g)$.}
\label{tab:direction}
\vspace{2pt}
\centering\small
\begin{tabular}{@{}lrrrr@{}}
\toprule
Configuration & $\cos(c,t)$ & covered, $\lambda$ (\%) & $\cos(c,g)$ & along generic (\%)\\
\midrule
GPT Image 1 & 0.85 & 44.1 & $-$0.03 & 0.1\\
GPT Image 2 & 0.66 & 50.7 & 0.09 & 0.9\\
Flare & 0.57 & 54.3 & 0.26 & 6.6\\
Sunburst & 0.72 & 55.6 & 0.47 & 22.0\\
Nano Banana 2 & 0.75 & 24.8 & 0.74 & 55.5\\
FLUX.2 Max & 0.86 & 64.6 & 0.73 & 52.8\\
\bottomrule
\end{tabular}
\end{table}

% Generated by paper/tmlr/build_assets.py; do not edit by hand.
% input reports/painter_tmlr_diagnostics_v1/analysis.json sha256 76792fa35f90e3d8fec1a75b0eec27d69cac14ed875fce8b64096b8e2525a4d4
% input reports/painter_specificity_v2/psv2-20260911/models.csv sha256 eed121f77626d8a37235b5bb19e370e3fbd7b6bae44b559d1b866e27e9cac41c
\begin{table}[t]
\caption{Further diagnostics in the 31 features. Repeat noise: half the squared difference between the two repeats' centered contrasts, relative to $H$. Corrected $H$: $\beta$ and $D$ after removing the finite-sample bias of $H$ (6.7\%); $\beta$, $Q$ and $D-1$ all scale by the same factor, so no ordering changes. Centroid proximity gain: mean reduction in squared distance from the generated mean to each painter's reference mean, from the generic to the named clause, split exactly into shared and between-name parts (Appendix~\ref{app:estimators}).}
\label{tab:robustness}
\vspace{2pt}
\centering\small
\begin{tabular}{@{}lrrrrrr@{}}
\toprule
 & Repeat noise & \multicolumn{2}{c}{Corrected $H$} & \multicolumn{3}{c}{Centroid proximity gain}\\
\cmidrule(lr){3-4}\cmidrule(l){5-7}
Configuration & $/H$ & $\beta$ & $D$ & total & shared & between-name\\
\midrule
GPT Image 1 & 2.02 & 1.006 & 1.614 & 17.490 & 17.844 & $-$0.354\\
GPT Image 2 & 0.96 & 1.071 & 1.242 & 5.094 & 5.159 & $-$0.065\\
Flare & 0.42 & 0.827 & 1.791 & 0.871 & 1.585 & $-$0.715\\
Sunburst & 0.58 & 0.883 & 1.688 & 4.500 & 4.997 & $-$0.498\\
Nano Banana 2 & 2.60 & 0.474 & 1.117 & 7.423 & 7.013 & 0.410\\
FLUX.2 Max & 1.67 & 0.504 & 0.787 & 10.529 & 10.176 & 0.353\\
\bottomrule
\end{tabular}
\end{table}

% Generated by paper/tmlr/build_assets.py; do not edit by hand.
% input reports/painter_specificity_v2/psv2-20260911/model_contrasts.csv sha256 0d0b722750e657b83658ccffbc7cb093d7b0ae577896904c45c02ababe31665b
\begin{table}[t]
\caption{All 15 prespecified pairwise error comparisons in the 31 features (positive values favor configuration B). Simultaneous: the prespecified 95\% intervals adjusted over the 21-endpoint family. Nominal: unadjusted paired-scene Student intervals. Bootstrap: the protocol's descriptive, unadjusted 95\% intervals from 5,000 paired scene resamples.}
\label{tab:pairs}
\vspace{2pt}
\centering\small\setlength{\tabcolsep}{4pt}
\begin{tabular}{@{}llrccc@{}}
\toprule
Configuration A & Configuration B & $D_A-D_B$ & simultaneous & nominal & bootstrap\\
\midrule
GPT Image 1 & GPT Image 2 & 0.346 & [$-$0.400, 1.093] & [$-$0.082, 0.775] & [$-$0.023, 0.728]\\
GPT Image 1 & Flare & $-$0.165 & [$-$1.027, 0.697] & [$-$0.661, 0.330] & [$-$0.579, 0.292]\\
GPT Image 1 & Sunburst & $-$0.069 & [$-$1.181, 1.043] & [$-$0.708, 0.570] & [$-$0.611, 0.518]\\
GPT Image 1 & Nano Banana 2 & 0.463 & [$-$0.602, 1.528] & [$-$0.149, 1.075] & [$-$0.049, 1.012]\\
GPT Image 1 & FLUX.2 Max & 0.771 & [$-$0.037, 1.580] & [0.307, 1.236] & [0.344, 1.176]\\
GPT Image 2 & Flare & $-$0.512 & [$-$1.150, 0.126] & [$-$0.878, $-$0.145] & [$-$0.823, $-$0.186]\\
GPT Image 2 & Sunburst & $-$0.415 & [$-$1.181, 0.350] & [$-$0.855, 0.024] & [$-$0.799, $-$0.044]\\
GPT Image 2 & Nano Banana 2 & 0.116 & [$-$0.922, 1.155] & [$-$0.480, 0.713] & [$-$0.443, 0.606]\\
GPT Image 2 & FLUX.2 Max & 0.425 & [$-$0.161, 1.011] & [0.088, 0.762] & [0.126, 0.718]\\
Flare & Sunburst & 0.096 & [$-$0.286, 0.478] & [$-$0.123, 0.316] & [$-$0.099, 0.283]\\
Flare & Nano Banana 2 & 0.628 & [$-$0.283, 1.540] & [0.105, 1.152] & [0.137, 1.057]\\
Flare & FLUX.2 Max & 0.937 & [0.256, 1.618] & [0.545, 1.328] & [0.589, 1.274]\\
Sunburst & Nano Banana 2 & 0.532 & [$-$0.436, 1.500] & [$-$0.024, 1.088] & [0.016, 1.006]\\
Sunburst & FLUX.2 Max & 0.840 & [0.058, 1.623] & [0.391, 1.290] & [0.422, 1.212]\\
Nano Banana 2 & FLUX.2 Max & 0.309 & [$-$0.847, 1.464] & [$-$0.355, 0.972] & [$-$0.220, 0.901]\\
\bottomrule
\end{tabular}
\end{table}

% Generated by paper/tmlr/build_assets.py; do not edit by hand.
% input reports/painter_specificity_review_v1/analysis.json sha256 2fabb78e15dfcc7b027a73170f6ad0c87af0196c754b4e6f45275269c2c961e2
% input data/manifests/painter_specificity_v2/psv2-20260911/analysis.json sha256 c1d21b16f840a10fd253c1d558743a09671acf3d4e48601ddcc96625a5a90be0
\begin{table}[t]
\caption{Parts of the error and contrast magnitude in the 31 features. Along and off-pattern: the prespecified exact split of $D$ into the error of the amplitude along the reference pattern, $(s_1-1)(s_2-1)$ with repeat-specific slopes $s_k$, and the cross-repeat product of the residuals orthogonal to that pattern, relative to $H$. $D=D_{\mathrm{agg}}+V_{\mathrm{scene}}$; $\beta/\sqrt{Q}$ is the alignment ratio. $D_{\mathrm{held}}$: error after multiplying every centered generated contrast by a nonnegative scalar fitted on the other 13 scenes; the fitted range covers the 14 folds. Rescaling acts on feature vectors only.}
\label{tab:calibration}
\vspace{2pt}
\centering\small\setlength{\tabcolsep}{3.5pt}
\begin{tabular}{@{}lrrrrrrcr@{}}
\toprule
 & & \multicolumn{2}{c}{Split of $D$} & \multicolumn{2}{c}{Scene split} & & &\\
\cmidrule(lr){3-4}\cmidrule(lr){5-6}
Configuration & $D$ & along & off-pattern & $D_{\mathrm{agg}}$ & $V_{\mathrm{scene}}$ & $\beta/\sqrt{Q}$ & fitted scalar & $D_{\mathrm{held}}$\\
\midrule
GPT Image 1 & 1.572 & 0.029 & 1.543 & 1.239 & 0.333 & 0.600 & 0.366--0.398 & 0.645\\
GPT Image 2 & 1.226 & $-$0.003 & 1.229 & 1.044 & 0.182 & 0.670 & 0.438--0.459 & \textbf{0.554}\\
Flare & 1.738 & 0.054 & 1.683 & 1.483 & 0.255 & 0.511 & 0.329--0.349 & 0.741\\
Sunburst & 1.641 & 0.044 & 1.598 & 1.337 & 0.305 & 0.545 & 0.346--0.373 & 0.706\\
Nano Banana 2 & 1.109 & 0.317 & 0.793 & 0.723 & 0.387 & 0.444 & 0.401--0.533 & 0.832\\
FLUX.2 Max & 0.801 & 0.304 & 0.497 & 0.761 & 0.040 & 0.546 & 0.598--0.694 & 0.714\\
\bottomrule
\end{tabular}
\end{table}

% Generated by paper/tmlr/build_assets.py; do not edit by hand.
% input reports/painter_tmlr_diagnostics_v3/analysis.json sha256 2aa36a87cd6d23ee909cdf94bd9f512d5743138f1d8a8f7c236d56f33fbceb57
% input reports/painter_specificity_review_v1/analysis.json sha256 2fabb78e15dfcc7b027a73170f6ad0c87af0196c754b4e6f45275269c2c961e2
\begin{table}[t]
\caption{Error $D$ of held-out genuine paintings sampled like the generated images (1,000 random half-splits; two works per painter for each of 14 pseudo-scenes, or 11 with content classes). With repl.: the recorded control, whose two pseudo-repeats can be the same work. Distinct works: the two pseudo-repeats are different paintings, as the two generated repeats are different images.}
\label{tab:genuine}
\vspace{2pt}
\centering\small\setlength{\tabcolsep}{3.5pt}
\begin{tabular}{@{}lrrcrcrc@{}}
\toprule
 & with repl. & \multicolumn{6}{c}{Distinct works}\\
\cmidrule(l){3-8}
 & 31 features & \multicolumn{2}{c}{31 features} & \multicolumn{2}{c}{CLIP} & \multicolumn{2}{c}{CSD}\\
Control & mean & mean & 95\% range & mean & 95\% range & mean & 95\% range\\
\midrule
pooled sampling, pooled target & 0.234 & 0.125 & [$-$0.886, 1.159] & 0.064 & [$-$0.162, 0.305] & 0.060 & [$-$0.209, 0.336]\\
class sampling, pooled target & 0.753 & 0.185 & [$-$1.030, 1.362] & 0.203 & [$-$0.106, 0.532] & 0.139 & [$-$0.177, 0.450]\\
class sampling, class targets & 0.731 & 0.328 & [$-$0.533, 1.166] & 0.216 & [$-$0.027, 0.470] & 0.193 & [$-$0.067, 0.465]\\
\bottomrule
\end{tabular}
\end{table}

% Generated by paper/tmlr/build_assets.py; do not edit by hand.
% input reports/painter_specificity_review_v1/analysis.json sha256 2fabb78e15dfcc7b027a73170f6ad0c87af0196c754b4e6f45275269c2c961e2
% input reports/painter_specificity_review_v3/analysis.json sha256 79fa5ca28900e852c072e2a52d35f3cb400ee053935d857729b4ad4aa3bc542c
% input reports/painter_learned_audit_v1/analysis.json sha256 06c77de4be133afc203dda122b2816d4f866e1389999f18bc94b3d9c668f28f7
\begin{table}[t]
\caption{Which painter differences carry the aligned response. Components: aligned amplitude along each singular component of the centered reference means (66.3\%, 20.6\% and 13.0\% of their spread). Omitting a painter recenters the other three and recomputes $\beta$. Monet--Sisley $\beta$: aligned amplitude for that pair alone, in the 31 features (full frame and central-square reference windows) and in each embedding.}
\label{tab:coverage}
\vspace{2pt}
\centering\footnotesize
\begin{tabular}{@{}l*{11}{r}@{}}
\toprule
 & \multicolumn{3}{c}{Component} & \multicolumn{4}{c}{$\beta$ after omitting} & \multicolumn{4}{c}{Monet--Sisley $\beta$}\\
\cmidrule(lr){2-4}\cmidrule(lr){5-8}\cmidrule(l){9-12}
Configuration & 1st & 2nd & 3rd & Mon. & Sis. & Pis. & C\'ez. & 31 & square & CLIP & CSD\\
\midrule
GPT Image 1 & 1.248 & 0.310 & 0.356 & 1.140 & 1.072 & 0.840 & 0.389 & 0.074 & 0.038 & 0.301 & 0.446\\
GPT Image 2 & 1.156 & 0.617 & 0.800 & 1.195 & 0.972 & 0.992 & 0.651 & 0.402 & 0.408 & 0.532 & 0.587\\
Flare & 1.022 & 0.211 & 0.383 & 1.034 & 0.734 & 0.802 & 0.222 & $-$0.011 & 0.012 & 0.339 & 0.451\\
Sunburst & 1.082 & 0.250 & 0.417 & 1.116 & 0.840 & 0.762 & 0.284 & $-$0.249 & $-$0.185 & 0.385 & 0.323\\
Nano Banana 2 & 0.513 & 0.371 & 0.195 & 0.569 & 0.446 & 0.388 & 0.276 & $-$0.044 & $-$0.017 & 0.318 & 0.400\\
FLUX.2 Max & 0.666 & 0.021 & 0.183 & 0.539 & 0.511 & 0.527 & 0.096 & 0.129 & 0.132 & 0.299 & 0.414\\
\bottomrule
\end{tabular}
\end{table}

% Generated by paper/tmlr/build_assets.py; do not edit by hand.
% input reports/painter_tmlr_diagnostics_v2/analysis.json sha256 bed891f767a98f4b281b6783b1021956f4d2c3ff80becbf6b2693ee45987ef1a
\begin{table}[t]
\caption{Monet--Sisley aligned amplitude with 95\% intervals from 5,000 scene resamples and 1,000 reference resamples (the latter shown for the 31 features; all six pairs and both interval types for every representation are in the supplementary analysis output).}
\label{tab:pair-intervals}
\vspace{2pt}
\centering\small\setlength{\tabcolsep}{4pt}
\begin{tabular}{@{}lrccrcrc@{}}
\toprule
 & \multicolumn{3}{c}{31 features} & \multicolumn{2}{c}{CLIP} & \multicolumn{2}{c}{CSD}\\
\cmidrule(lr){2-4}\cmidrule(lr){5-6}\cmidrule(l){7-8}
Configuration & $\beta$ & scenes & references & $\beta$ & scenes & $\beta$ & scenes\\
\midrule
GPT Image 1 & 0.074 & [$-$0.166, 0.310] & [$-$0.104, 0.193] & 0.301 & [0.239, 0.363] & 0.446 & [0.355, 0.539]\\
GPT Image 2 & 0.402 & [0.238, 0.564] & [0.154, 0.617] & 0.532 & [0.456, 0.614] & 0.587 & [0.511, 0.669]\\
Flare & $-$0.011 & [$-$0.088, 0.052] & [$-$0.159, 0.204] & 0.339 & [0.286, 0.394] & 0.451 & [0.391, 0.509]\\
Sunburst & $-$0.249 & [$-$0.369, $-$0.140] & [$-$0.443, 0.122] & 0.385 & [0.311, 0.465] & 0.323 & [0.248, 0.402]\\
Nano Banana 2 & $-$0.044 & [$-$0.285, 0.208] & [$-$0.235, 0.230] & 0.318 & [0.259, 0.377] & 0.400 & [0.333, 0.478]\\
FLUX.2 Max & 0.129 & [$-$0.069, 0.323] & [0.001, 0.232] & 0.299 & [0.201, 0.394] & 0.414 & [0.334, 0.496]\\
\bottomrule
\end{tabular}
\end{table}

% Generated by paper/tmlr/build_assets.py; do not edit by hand.
% input data/manifests/painter_specificity_v2/psv2-20260911/analysis.json sha256 c1d21b16f840a10fd253c1d558743a09671acf3d4e48601ddcc96625a5a90be0
\begin{table}[t]
\caption{Prespecified feature-family sensitivity of the agreement scores: $\beta$ and $D$ computed within each family (with its own $H$) and without the texture family.}
\label{tab:family-agreement}
\vspace{2pt}
\centering\small\setlength{\tabcolsep}{4pt}
\begin{tabular}{@{}lrrrrrrrr@{}}
\toprule
 & \multicolumn{4}{c}{Aligned amplitude $\beta$} & \multicolumn{4}{c}{Error $D$}\\
\cmidrule(lr){2-5}\cmidrule(l){6-9}
Configuration & color & spatial & texture & no texture & color & spatial & texture & no texture\\
\midrule
GPT Image 1 & 0.905 & 0.835 & 1.037 & 0.876 & 0.889 & 1.552 & 2.220 & 1.163\\
GPT Image 2 & 1.061 & 1.112 & 0.866 & 1.082 & 1.029 & 1.433 & 1.274 & 1.196\\
Flare & 0.810 & 1.136 & 0.497 & 0.945 & 1.193 & 2.129 & 1.987 & 1.580\\
Sunburst & 0.943 & 1.070 & 0.553 & 0.995 & 1.256 & 2.097 & 1.701 & 1.604\\
Nano Banana 2 & 0.610 & 0.487 & 0.257 & 0.559 & 0.664 & 2.084 & 0.886 & 1.251\\
FLUX.2 Max & 0.465 & 0.394 & 0.524 & 0.436 & 0.605 & 0.917 & 0.906 & 0.734\\
\bottomrule
\end{tabular}
\end{table}

% Generated by paper/tmlr/build_assets.py; do not edit by hand.
% input reports/painter_specificity_review_v1/analysis.json sha256 2fabb78e15dfcc7b027a73170f6ad0c87af0196c754b4e6f45275269c2c961e2
% input reports/painter_reference_quality_v1/analysis.json sha256 089f636335129d6377328088ccc9fdc763488aba2a88463939322326707052b4
% input reports/painter_specificity_v2/psv2-20260911/models.csv sha256 eed121f77626d8a37235b5bb19e370e3fbd7b6bae44b559d1b866e27e9cac41c
% input data/manifests/painter_specificity_v2/psv2-20260911/analysis.json sha256 c1d21b16f840a10fd253c1d558743a09671acf3d4e48601ddcc96625a5a90be0
\begin{table}[t]
\caption{Error $D$ (31 features) under alternative scene sets, windows, targets, weightings and source corrections. Deletions: range over the 14 single-scene deletions. Square: central-square measurement windows. Class target: title-derived water, built and land reference means for the 11 matching scenes (pooled: the same 11 scenes against the pooled target). Equal: each feature family has equal weight; covar.: shrunk inverse covariance of the development panel. Cropped: AI-audited painting regions with the original scaling; refit: also refitting the development scaling.}
\label{tab:sensitivity}
\vspace{2pt}
\centering\small\setlength{\tabcolsep}{3.5pt}
\begin{tabular}{@{}lrcrrrrrrr@{}}
\toprule
 & & & & \multicolumn{2}{c}{11 scenes} & \multicolumn{2}{c}{Weighting} & \multicolumn{2}{c}{Source correction}\\
\cmidrule(lr){5-6}\cmidrule(lr){7-8}\cmidrule(l){9-10}
Configuration & Primary & deletions & square & pooled & class & equal & covar. & cropped & refit\\
\midrule
GPT Image 1 & 1.572 & 1.469--1.694 & 1.436 & 1.557 & 1.366 & 1.551 & 2.103 & 1.561 & 1.631\\
GPT Image 2 & 1.226 & 1.175--1.285 & 1.128 & 1.099 & 1.130 & 1.243 & 1.511 & 1.171 & 1.259\\
Flare & 1.738 & 1.674--1.799 & 1.649 & 1.726 & 1.539 & 1.765 & 1.636 & 1.680 & 1.737\\
Sunburst & 1.641 & 1.568--1.729 & 1.534 & 1.633 & 1.525 & 1.680 & 1.617 & 1.536 & 1.577\\
Nano Banana 2 & 1.109 & 0.948--1.205 & 1.024 & 0.987 & 0.948 & 1.203 & 1.385 & 1.076 & 1.121\\
FLUX.2 Max & 0.801 & 0.754--0.852 & 0.765 & 0.750 & 0.786 & 0.808 & 0.928 & 0.832 & 0.872\\
\bottomrule
\end{tabular}
\end{table}


\begin{figure}[ht]
\centering
\includegraphics[width=\linewidth]{fig_pairs.pdf}
\caption{Aligned amplitude and error for each of the six painter pairs (M: Monet, S: Sisley,
P: Pissarro, C: C\'ezanne), using all 31 features and normalized by the squared reference
distance of each pair. An amplitude of 1 matches the reference, negative values point the
opposite way, and an error of 1 corresponds to no distinction. Negative errors can arise from
repeat correction.}
\label{fig:pairs}
\end{figure}

\paragraph{Full decomposition.} Table~\ref{tab:decomposition-full} gives all components; against the
artist-free baseline the shared component is $N_{\mathrm{free}}=G+N+I$, with $G$ the squared
generic-minus-free shift and $I$ a signed cross term (Appendix~\ref{app:estimators}). Against
the artist-free baseline, 82.5--95.7\% of the squared change is shared; this includes the
painting instruction itself, and computing both components within scenes gives 88.6--97.2\%.
Against the generic baseline the within-scene shares are 52.8--93.6\%. Table~\ref{tab:robustness}
reports centroid proximity in squared standardized units.

\paragraph{Feature families.} Table~\ref{tab:families} gives the shared fraction by family, with the
mean faithful and exact-differences values. The faithful benchmark is higher for texture than for
color (+2.1 points), the exact-differences benchmark lower ($-$3.2 points; both computed before
rounding), so neither benchmark explains the texture--color difference alone. Reweighting the 31
features gives shared fractions of 66.1--89.8\% (equal family weights) and 65.8--86.9\%
(development covariance), and deleting any single feature keeps the fraction within 64.8--89.6\%. Scene resamples in
which a family's denominator was not positive were dropped: 32 of 5,000 for the spatial and 4 for
the texture mean, and 35 for the paired family differences.

\paragraph{Genuine-painting controls.} For each of 1,000 random splits, the works in each
painter's collection (or each painter and content class) are divided into two disjoint halves;
the smaller half estimates the target, and the other supplies two distinct works per painter for
each of 14 pseudo-scenes (11 with content classes), which are scored with the same error $D$ in
each representation (Table~\ref{tab:genuine}). An earlier version of this control drew the two
works with replacement, so that both pseudo-repeats could be the same painting; the product of a
work with itself adds its within-painter variance, which inflated that version's means (0.234,
0.753 and 0.731, first column). With distinct works the pseudo-repeats are unbiased for the
painter's mean, so the pooled control's expected value is the noise of the target alone: for a
target estimated from about half of each collection, the term $\frac34\sum_a\operatorname{tr}
(\Sigma_a)/n_a$ of Appendix~\ref{app:estimators} is about twice the 6.7\% of $H$ found for the full
collections, close to the observed 0.125. Content-class sampling adds the difference between class
and pooled means (0.185), and class-specific targets rest on fewer works (0.328). The generated
errors are scored against the full collections, whose target noise is the 6.7\% removed by the
correction of $H$ (Table~\ref{tab:robustness}). These controls therefore describe what genuine
paintings score under this design; they do not bound what a generator could reach.

\paragraph{Feature separability.} Classifying each of the 221 development works by its nearest
reference mean in the 31 standardized features gives 49.8\% macro accuracy, with 39.6\% of Monet
and 30.6\% of Sisley works correct; the embeddings give 79.8\% (CLIP) and 79.6\% (CSD).

\paragraph{Painter pairs and coverage.} Table~\ref{tab:coverage} shows that the aggregate
aligned response rests mostly on the first reference component, which accounts for 66.3\% of
the reference spread, and that omitting C\'ezanne removes much of it. Stronger pair alignment
does not ensure lower pair error: GPT Image 1 has weaker Monet--Sisley alignment than GPT Image
2 but lower error for that pair (Figure~\ref{fig:pairs}). For Flare, Sunburst and Nano Banana
2, swapping the Monet and Sisley reference labels lowers the overall error, as it must when the
pair alignment is negative.

\paragraph{Scene aggregation and rescaling.} Nano Banana 2 has the largest scene variation and
FLUX.2 Max the smallest, which explains their reversal between $D$ and $D_{\mathrm{agg}}$. GPT
Image 2's differences are about twice the reference size in squared size (1.49 times in norm), so
shrinking them to about 0.45 of their size removes more off-pattern difference than aligned
response. Deleting any scene and
refitting the whole procedure preserves GPT Image 2's advantage after rescaling.

\paragraph{Distributions.} Table~\ref{tab:diversity} gives the prespecified distribution
diagnostics. The 28 images named for a painter span 14 scenes, yet their total variance is 0.25--0.88
of that of the painter's reference works. The named images are closer than the generic images to the
painter's works in energy distance in 21 of 24 configuration--painter combinations.
Figure~\ref{fig:projection} shows each configuration's scene-averaged centered named means next to
the reference differences on the first two reference axes, which carry 87.0\% of $H$.

% Generated by paper/tmlr/build_assets.py; do not edit by hand.
% input data/manifests/painter_specificity_v2/psv2-20260911/analysis.json sha256 c1d21b16f840a10fd253c1d558743a09671acf3d4e48601ddcc96625a5a90be0
\begin{table}[t]
\caption{Prespecified distribution diagnostics in the 31 features. Trace ratio: total variance of the 28 images named for a painter (14 scenes, two repeats) relative to that of the painter's reference works. Energy distance: between those images and the painter's reference works, averaged over painters, for the named and for the generic images; named closer: number of painters (of four) for which the named images are closer.}
\label{tab:diversity}
\vspace{2pt}
\centering\small
\begin{tabular}{@{}lrrrrrrr@{}}
\toprule
 & \multicolumn{4}{c}{Trace ratio} & \multicolumn{3}{c}{Energy distance}\\
\cmidrule(lr){2-5}\cmidrule(l){6-8}
Configuration & Mon. & Sis. & Pis. & C\'ez. & named & generic & named closer\\
\midrule
GPT Image 1 & 0.48 & 0.69 & 0.60 & 0.60 & 2.64 & 4.94 & 4\\
GPT Image 2 & 0.28 & 0.31 & 0.30 & 0.30 & 2.20 & 2.76 & 4\\
Flare & 0.27 & 0.37 & 0.32 & 0.27 & 2.11 & 2.01 & 2\\
Sunburst & 0.25 & 0.27 & 0.27 & 0.32 & 1.81 & 2.36 & 3\\
Nano Banana 2 & 0.83 & 0.71 & 0.57 & 0.88 & 2.50 & 3.43 & 4\\
FLUX.2 Max & 0.34 & 0.42 & 0.57 & 0.62 & 1.35 & 2.37 & 4\\
\bottomrule
\end{tabular}
\end{table}


\begin{figure}[ht]
\centering
\includegraphics[width=\linewidth]{fig_projection.pdf}
\caption{Reference differences $r_a$ (squares) and each configuration's scene- and
repeat-averaged centered named means (circles) on the first two principal axes of the 31-feature
reference differences (M: Monet, S: Sisley, P: Pissarro, C: C\'ezanne; gray lines join each
painter's two points). Point values; the remaining axes are not shown.}
\label{fig:projection}
\end{figure}

\paragraph{Reference sensitivity.} The content-matched target uses the title-derived water,
built and land classes for the 11 non-mixed scenes; works in the route class enter the pooled
target but no class target, because no scene belongs to that class. The source audit recorded removable
frames, margins, labels and calibration strips; 43 uncertain boundaries were left unchanged. A
visual content audit of the same reproductions disagreed with the title-derived class for 230
works; replacing the title-derived with visual labels also leaves FLUX.2 Max with the lowest
error. All works and features are retained under every correction.

\FloatBarrier
\section{Feature families and sensitivity}\label{app:sensitivity}

\paragraph{Feature families.} Averaged over configurations, the shared fraction is 74.3\% in color,
71.5\% in spatial and 68.2\% in texture features, but the exact-differences benchmark is also lower
for texture (74.7\%, against 77.8\% for color and 79.1\% for spatial), and every family lies 3.6--7.6
points below its benchmark (Table~\ref{tab:families}). The texture--color difference of
$-$6.1 points [$-$13.5, $-$2.0] is therefore partly expected from the reference differences, and the
texture--spatial ordering is not resolved (texture lower in 65.2\% of paired scene resamples). GPT
Image 1's texture fraction, 48.6\%, the only value below 50\%, comes from oversized texture
differences ($B/H=$ 3.11; exact-differences value 74.6\%). An earlier collection of 2,000 SD-Turbo
images with a different template and seed design shares 64.2\% of the change across all features
and 36.6\% in texture, against faithful and exact-differences values of 92.5\% and 67.4\% overall
(Appendix~\ref{app:sdturbo}); it reuses the references and serves as a retrospective check.

\paragraph{Reference sources.} The shared fraction does not depend on the reference collection; its
benchmarks and the agreement scores do. AI assistants audited all 870 reference and development
reproductions and cropped peripheral artifacts such as frames and calibration strips from 131 of
them; the crops were not verified by a human. With the cropped regions and the original feature
scaling, FLUX.2 Max's error becomes 0.832, all six aligned amplitudes remain above zero, and of the
pairwise differences only Flare--FLUX.2 Max remains resolved. After also refitting the feature
scaling, FLUX.2 Max's error is 0.872, the amplitudes remain above zero, and Flare--FLUX.2 Max and GPT
Image 1--FLUX.2 Max are resolved.

\paragraph{Other checks.} FLUX.2 Max keeps the lowest 31-feature error estimate under every
single-scene deletion, square measurement windows, content-matched reference targets, both
alternative feature weightings and both source corrections (Table~\ref{tab:sensitivity}),
but in texture features alone Nano Banana 2 is lower, 0.886 against 0.906 (Appendix
Table~\ref{tab:family-agreement}). Correcting the reference spread $H$ for its finite-sample bias
(6.7\%) scales $\beta$, $Q$ and $D-1$ by the same factor: the aligned amplitudes become
0.474--1.071, FLUX.2 Max's error 0.787, no ordering changes, and the exact-differences fractions
rise to 59.3--85.1\% (Table~\ref{tab:robustness}). The named images of every configuration
are less varied than the painter's reference works (trace ratios 0.25--0.88), and they are closer
to the painter's works than the generic images in 21 of 24 configuration--painter combinations
(Table~\ref{tab:diversity}).

\FloatBarrier
\section{Learned embeddings}\label{app:learned}

\paragraph{Extraction.} We used the public checkpoints
\texttt{openai/\allowbreak clip-\allowbreak vit-\allowbreak large-\allowbreak patch14} and
\texttt{tomg-\allowbreak group-\allowbreak umd/\allowbreak CSD-\allowbreak ViT-\allowbreak L}. CSD uses its released style projection of the CLIP visual
backbone; its release notes a discrepancy with the published benchmark numbers, so we
characterize the released checkpoint without claiming to reproduce those numbers. Every image
is orientation- and color-corrected as above, resized bicubically to a 224-pixel short side,
center-cropped to $224\times224$ and normalized with the CLIP channel statistics. Embeddings
are unit-normalized 768-dimensional vectors; we use their Euclidean and cosine geometry without
scaling or dimension selection. Backend and batch-size checks on fixed test cases agreed to
within $1.3\times10^{-6}$ per coordinate.

\paragraph{Prototype gain.} With unnormalized prototypes $\mu_a$ (means of unit reference
embeddings), $g_a^\top\mu_b$ is the average image-to-reference cosine similarity. Writing
$g_a=\bar g+d_a$ and $\mu_a=\bar\mu+r_a$,
\begin{equation}
\frac14\sum_a g_a^\top\mu_a-g_{\mathrm g}^\top\bar\mu
=(\bar g-g_{\mathrm g})^\top\bar\mu+\frac14\sum_a d_a^\top r_a
=(\bar g-g_{\mathrm g})^\top\bar\mu+\frac{H\widehat\beta}{4},
\end{equation}
because $\sum_a r_a=\sum_a d_a=0$. The first term cannot depend on which name produced which
images; permuting the painter labels of the generated images changes only $\widehat\beta$. The
faithful value sets $g_a=\mu_a$, giving the shared fraction $(\bar\mu-g_{\mathrm
g})^\top\bar\mu/[(\bar\mu-g_{\mathrm g})^\top\bar\mu+H/4]$. Replacing $\mu_a$ by the unit
prototype $u_a=\mu_a/\|\mu_a\|$ keeps the score linear in the image embedding, so the same
identity holds with $u_a$ in place of $\mu_a$ (Table~\ref{tab:learned-settings}).

% Generated by paper/tmlr/build_assets.py; do not edit by hand.
% input reports/painter_tmlr_diagnostics_v1/analysis.json sha256 76792fa35f90e3d8fec1a75b0eec27d69cac14ed875fce8b64096b8e2525a4d4
% input reports/painter_tmlr_diagnostics_v5/analysis.json sha256 5747902e4fae5ff111a52b5f2b8da112bc26c611555e82314b278de953981f6b
\begin{table}[t]
\caption{Which configuration is best on each readout when the 14 scenes are resampled with replacement (5,000 paired resamples; the same scenes for every configuration). The percentages describe dependence on the authored scenes. Ties, which occur for recognition, are credited to the configuration listed first. Lower rows: alignment ratio and $D_{\mathrm{held}}$ (separate resamples).}
\label{tab:stability}
\vspace{2pt}
\centering\small
\begin{tabular}{@{}llrlr@{}}
\toprule
Readout & Most often best & \% & Next & \%\\
\midrule
Proximity gain, CLIP & Nano Banana 2 & 85.8 & Sunburst & 9.5\\
Proximity gain, CSD & FLUX.2 Max & 33.6 & GPT Image 1 & 26.2\\
Recognition, CLIP & GPT Image 2 & 98.7 & GPT Image 1 & 0.8\\
Recognition, CSD & GPT Image 2 & 80.4 & GPT Image 1 & 19.6\\
Error $D$, 31 features & FLUX.2 Max & 84.7 & Nano Banana 2 & 15.3\\
Error $D$, CLIP & GPT Image 1 & 99.3 & GPT Image 2 & 0.6\\
Error $D$, CSD & FLUX.2 Max & 47.2 & GPT Image 1 & 26.8\\
\midrule
Alignment ratio, 31 features & GPT Image 2 & 98.0 & Nano Banana 2 & 0.7\\
Alignment ratio, CLIP & GPT Image 2 & 99.7 & GPT Image 1 & 0.3\\
Alignment ratio, CSD & GPT Image 2 & 100.0 & FLUX.2 Max & 0.0\\
$D_{\mathrm{held}}$, 31 features & GPT Image 2 & 99.0 & GPT Image 1 & 0.6\\
$D_{\mathrm{held}}$, CLIP & GPT Image 2 & 99.8 & GPT Image 1 & 0.2\\
$D_{\mathrm{held}}$, CSD & GPT Image 2 & 100.0 & FLUX.2 Max & 0.0\\
\bottomrule
\end{tabular}
\end{table}

% Generated by paper/tmlr/build_assets.py; do not edit by hand.
% input reports/painter_learned_audit_v1/analysis.json sha256 06c77de4be133afc203dda122b2816d4f866e1389999f18bc94b3d9c668f28f7
% input reports/painter_tmlr_diagnostics_v3/analysis.json sha256 2aa36a87cd6d23ee909cdf94bd9f512d5743138f1d8a8f7c236d56f33fbceb57
\begin{table}[t]
\caption{Share (\%) of the proximity gain supplied by the shared term under alternative reference settings. Prim.: full reproductions of the primary reference collection (Table~\ref{tab:learned}). Crop: reference reproductions replaced by their audited painting regions. Dev.: the development panel as the reference target. Both: audited regions of the development panel. Unit: primary prototypes normalized to unit length, as in recognition.}
\label{tab:learned-settings}
\vspace{2pt}
\centering\small\setlength{\tabcolsep}{3.5pt}
\begin{tabular}{@{}lrrrrrrrrrr@{}}
\toprule
 & \multicolumn{5}{c}{CLIP} & \multicolumn{5}{c}{CSD}\\
\cmidrule(lr){2-6}\cmidrule(l){7-11}
Configuration & prim. & crop & dev. & both & unit & prim. & crop & dev. & both & unit\\
\midrule
GPT Image 1 & 74.5 & 74.5 & 74.3 & 74.4 & 74.6 & 72.7 & 72.9 & 73.4 & 73.0 & 72.8\\
GPT Image 2 & 73.2 & 73.0 & 72.8 & 72.7 & 73.4 & 54.2 & 54.7 & 53.5 & 53.2 & 54.7\\
Flare & 78.8 & 78.7 & 78.9 & 79.0 & 79.0 & 65.0 & 65.5 & 65.1 & 64.9 & 65.4\\
Sunburst & 77.7 & 77.7 & 78.0 & 78.1 & 78.0 & 67.4 & 67.6 & 68.0 & 67.8 & 67.7\\
Nano Banana 2 & 83.8 & 83.9 & 83.5 & 83.7 & 83.9 & 79.7 & 79.9 & 81.3 & 81.0 & 79.6\\
FLUX.2 Max & 81.7 & 81.7 & 81.5 & 81.9 & 81.9 & 77.6 & 77.7 & 78.7 & 78.5 & 77.7\\
\bottomrule
\end{tabular}
\end{table}

% Generated by paper/tmlr/build_assets.py; do not edit by hand.
% input reports/painter_tmlr_diagnostics_v5/analysis.json sha256 5747902e4fae5ff111a52b5f2b8da112bc26c611555e82314b278de953981f6b
\begin{table}[t]
\caption{Error parts and repeat dependence in each representation. $\beta$: unadjusted paired-scene Student 95\% intervals. Split of $D$: the prespecified parts along and off the reference pattern (Table~\ref{tab:calibration}) and the off-pattern share. Repeat noise: half the squared difference between the repeats' centered contrasts, relative to $H$. $\rho$ at $D=1$: common repeat correlation at which an error above 1 would fall to 1 if the repeats shared that fraction of their noise (Appendix~\ref{app:provenance}); errors below 1 would only fall further.}
\label{tab:embedding-calibration}
\vspace{2pt}
\centering\small\setlength{\tabcolsep}{4pt}
\begin{tabular}{@{}lcrrrrrr@{}}
\toprule
 & $\beta$ & & \multicolumn{3}{c}{Split of $D$} & Repeat &\\
\cmidrule(lr){4-6}
Configuration & scenes & $D$ & along & off-pattern & off (\%) & noise $/H$ & $\rho$ at $D=1$\\
\midrule
\multicolumn{8}{@{}l}{\textit{31 features}}\\
GPT Image 1 & [0.806, 1.070] & 1.572 & 0.029 & 1.543 & 98.2 & 2.02 & 0.22\\
GPT Image 2 & [0.938, 1.059] & 1.226 & $-$0.003 & 1.229 & 100.2 & 0.96 & 0.19\\
Flare & [0.719, 0.824] & 1.738 & 0.054 & 1.683 & 96.9 & 0.42 & 0.64\\
Sunburst & [0.741, 0.907] & 1.641 & 0.044 & 1.598 & 97.3 & 0.58 & 0.52\\
Nano Banana 2 & [0.335, 0.549] & 1.109 & 0.317 & 0.793 & 71.5 & 2.60 & 0.04\\
FLUX.2 Max & [0.337, 0.603] & 0.801 & 0.304 & 0.497 & 62.0 & 1.67 & ---\\
\midrule
\multicolumn{8}{@{}l}{\textit{CLIP}}\\
GPT Image 1 & [0.468, 0.585] & 0.847 & 0.231 & 0.616 & 72.8 & 1.20 & ---\\
GPT Image 2 & [0.716, 0.825] & 1.061 & 0.057 & 1.004 & 94.6 & 1.13 & 0.05\\
Flare & [0.511, 0.606] & 1.057 & 0.201 & 0.856 & 81.0 & 0.80 & 0.07\\
Sunburst & [0.594, 0.709] & 1.136 & 0.129 & 1.007 & 88.6 & 0.83 & 0.14\\
Nano Banana 2 & [0.472, 0.574] & 1.078 & 0.230 & 0.848 & 78.7 & 1.59 & 0.05\\
FLUX.2 Max & [0.385, 0.490] & 1.058 & 0.316 & 0.743 & 70.2 & 1.78 & 0.03\\
\midrule
\multicolumn{8}{@{}l}{\textit{CSD}}\\
GPT Image 1 & [0.654, 0.803] & 0.735 & 0.085 & 0.650 & 88.4 & 0.73 & ---\\
GPT Image 2 & [0.889, 0.982] & 0.741 & 0.005 & 0.736 & 99.4 & 0.54 & ---\\
Flare & [0.768, 0.841] & 0.819 & 0.041 & 0.778 & 95.0 & 0.36 & ---\\
Sunburst & [0.813, 0.906] & 0.944 & 0.024 & 0.920 & 97.4 & 0.42 & ---\\
Nano Banana 2 & [0.463, 0.588] & 0.999 & 0.231 & 0.768 & 76.9 & 1.04 & ---\\
FLUX.2 Max & [0.540, 0.656] & 0.714 & 0.162 & 0.551 & 77.3 & 1.01 & ---\\
\bottomrule
\end{tabular}
\end{table}

% Generated by paper/tmlr/build_assets.py; do not edit by hand.
% input reports/painter_tmlr_diagnostics_v1/analysis.json sha256 76792fa35f90e3d8fec1a75b0eec27d69cac14ed875fce8b64096b8e2525a4d4
\begin{table}[t]
\caption{Shared fraction $N/(N+B)$ (\%) of the squared change added by the names, computed in each embedding, with its faithful-imitation value.}
\label{tab:learned-squared}
\vspace{2pt}
\centering\small
\begin{tabular}{@{}lrrrr@{}}
\toprule
 & \multicolumn{2}{c}{CLIP} & \multicolumn{2}{c}{CSD}\\
\cmidrule(lr){2-3}\cmidrule(l){4-5}
Configuration & observed & faithful & observed & faithful\\
\midrule
GPT Image 1 & 76.1 & 86.1 & 76.5 & 89.2\\
GPT Image 2 & 76.6 & 87.1 & 68.8 & 82.4\\
Flare & 81.7 & 88.7 & 74.6 & 86.4\\
Sunburst & 77.6 & 88.9 & 69.9 & 86.9\\
Nano Banana 2 & 82.3 & 91.2 & 77.4 & 89.7\\
FLUX.2 Max & 80.8 & 88.5 & 75.1 & 87.7\\
\bottomrule
\end{tabular}
\end{table}

% Generated by paper/tmlr/build_assets.py; do not edit by hand.
% input reports/painter_tmlr_diagnostics_v2/analysis.json sha256 bed891f767a98f4b281b6783b1021956f4d2c3ff80becbf6b2693ee45987ef1a
\begin{table}[t]
\caption{Shared part (\%) of the proximity gain for each prompted painter separately: $(\bar g-g_{\mathrm g})^\top\mu_a$ divided by $(g_a-g_{\mathrm g})^\top\mu_a$. Values above 100\% mean that the painter-specific term is negative for that painter.}
\label{tab:per-painter}
\vspace{2pt}
\centering\small\setlength{\tabcolsep}{4pt}
\begin{tabular}{@{}lrrrrrrrr@{}}
\toprule
 & \multicolumn{4}{c}{CLIP} & \multicolumn{4}{c}{CSD}\\
\cmidrule(lr){2-5}\cmidrule(l){6-9}
Configuration & Mon. & Sis. & Pis. & C\'ez. & Mon. & Sis. & Pis. & C\'ez.\\
\midrule
GPT Image 1 & 75.4 & 85.7 & 81.2 & 59.5 & 73.8 & 79.6 & 73.1 & 62.0\\
GPT Image 2 & 78.1 & 85.5 & 91.9 & 51.6 & 68.8 & 62.1 & 64.1 & 34.6\\
Flare & 84.7 & 93.3 & 102.9 & 55.1 & 74.1 & 68.1 & 85.8 & 47.5\\
Sunburst & 90.2 & 83.7 & 93.5 & 57.2 & 75.3 & 70.4 & 78.4 & 53.4\\
Nano Banana 2 & 87.6 & 90.0 & 87.6 & 73.2 & 79.1 & 79.1 & 73.4 & 85.8\\
FLUX.2 Max & 80.5 & 105.3 & 97.0 & 57.6 & 70.6 & 86.4 & 96.9 & 64.0\\
\bottomrule
\end{tabular}
\end{table}


\paragraph{Error parts and repeat dependence.} Table~\ref{tab:embedding-calibration} gives the split
of $D$, Student intervals for $\beta$ and the repeat noise in all three representations. In the
embeddings the off-pattern part also dominates, and the amplitude along the pattern contributes
more than in the features.

\paragraph{Crops.} Generated images are square, so the encoders' center crop sees nearly the whole
image, whereas non-square reference reproductions lose their periphery. For the 31 features the
comparable square-window sensitivity leaves FLUX.2 Max with the lowest error (Appendix
Table~\ref{tab:sensitivity}).

\paragraph{Scope.} CSD is built on the CLIP backbone and both encoders see the same images, so
their agreement is weaker evidence than agreement between independent representations would be. All
four painters appear in CSD's style-tag vocabulary, and the overlap between either encoder's
training data and our references is unknown. As a check on the reference panels, the primary
prototypes classify the 221 development works with 79.8\% (CLIP) and 79.6\% (CSD) macro
accuracy.

\FloatBarrier
\section{Held-scene prompt-name recognition}\label{app:transfer}

% Generated by paper/tmlr/build_assets.py; do not edit by hand.
% input reports/painter_prototype_transfer_v1/analysis.json sha256 e1d1fbb924feb797f67e0907677511ed7c745335a4f2827e99d2f286e6be5448
\begin{table}[t]
\caption{Held-scene recognition of the prompted painter (\%, 112 named images per configuration, chance 25\%). Ref.: nearest reference prototype. Shift: the same after adding one fixed translation that aligns the generated and reference mixture means, fitted on the other 13 scenes without painter labels. Gen.: nearest generated-image prototype fitted with painter labels on the other 13 scenes. Shift and Gen.\ use more information than Ref., and Gen.\ more than Shift.}
\label{tab:transfer}
\vspace{2pt}
\centering\small
\begin{tabular}{@{}lrrrrrr@{}}
\toprule
 & \multicolumn{3}{c}{CLIP} & \multicolumn{3}{c}{CSD}\\
\cmidrule(lr){2-4}\cmidrule(l){5-7}
Configuration & Ref. & Shift & Gen. & Ref. & Shift & Gen.\\
\midrule
GPT Image 1 & 60.7 & 55.4 & 80.4 & 72.3 & 73.2 & 84.8\\
GPT Image 2 & 68.8 & 66.1 & 75.0 & 75.9 & 76.8 & 94.6\\
Flare & 59.8 & 59.8 & 77.7 & 62.5 & 74.1 & 92.9\\
Sunburst & 58.9 & 62.5 & 69.6 & 61.6 & 70.5 & 90.2\\
Nano Banana 2 & 51.8 & 59.8 & 63.4 & 50.9 & 61.6 & 75.0\\
FLUX.2 Max & 41.1 & 53.6 & 71.4 & 39.3 & 66.1 & 75.0\\
\midrule
Mean change vs.\ Ref.\ (pp) & & +2.7 & +16.1 & & +10.0 & +25.0\\
\bottomrule
\end{tabular}
\end{table}


For each configuration and held-out scene $s$, let $x$ be a unit embedding of a named image
from scene $s$ and $u_a=\mu_a/\|\mu_a\|$ the normalized reference prototypes. The reference rule
(Ref.) is $\arg\max_a x^\top u_a$. The shift rule computes $t_{-s}=\frac14\sum_a\mu_a-\bar
g_{-s}$, where $\bar g_{-s}$ is the mean of the 104 named embeddings from the other 13 scenes,
and predicts $\arg\max_a(x+t_{-s})^\top u_a$. The translation adds $t_{-s}^\top u_a$ to painter
$a$'s score; these shifts differ between painters, which is why recognition can change. Because the
same vector is added to every named image, the centered contrasts $d_{sak}$, and hence $\beta$, $Q$
and $D$, are unchanged. The generated-prototype rule (Gen.) uses the normalized mean of the 26
labeled training embeddings for each painter. Held-scene repeats never enter the fit, and the
artist-free and generic images are not used. The question was chosen after the reference-rule
confusions were known; the rules were fixed before the translated outcomes were computed. Folds
overlap and scenes are authored, so these are descriptive out-of-fold results. The mean changes
from the reference rule are +2.7 (CLIP) and +10.0 points (CSD) for the shift rule; for FLUX.2 Max in
CSD the shift corrects 39 and harms 9 of 112 decisions. The generated-prototype rule outperforms
both reference-based rules in all 12 configuration--encoder combinations, so the prompted names
are more separable within each generator's own outputs than with respect to the historical
references. Across the four
combinations of reference source (full reproductions or audited painting regions) and reference
target (primary collection or development panel), the mean translation gains are 2.2--4.0 points
(CLIP) and 9.7--10.1 points (CSD).

\FloatBarrier
\section{The SD-Turbo collection}\label{app:sdturbo}

The collection was generated on 5 September 2026 with \texttt{stabilityai/sd-turbo} at
$512\times512$ pixels, one denoising step and guidance scale zero, for an earlier, unpublished
distance analysis by the same project. It crosses 16 scenes, 25 repeat blocks and five
conditions (the baseline and four names), 400 images per condition. Within each scene and
block, all five conditions use the same seed. Its image hashes are disjoint from the main
collection, and its 16 scene descriptions are shorter, class-level texts, different from the 14
main descriptions (for example, ``a riverbank landscape, with water organizing the
composition''), and the baseline prompt already requests an oil painting, with the names inserted
as ``by [painter]''.

For block vectors $v_b$, $b=1,\dots,K$ with $K=25$, the estimator
$U(v)=[K(K-1)]^{-1}\sum_{b\ne b'}\ip{v_b}{v_{b'}}$ replaces the two-repeat product. Averaging
scenes within each block, let $\delta_{ba}$ be the named-minus-baseline shift, $c_b$ its mean
over names and $e_{ba}=\delta_{ba}-c_b$. The shared component is $4U(c)$ and the between-name
component is $\sum_aU(e_a)$. Because products are formed only between different blocks,
correlation among the five seed-matched conditions within a block is allowed. The reference
means and scaling are those of the main experiment. Deleting each block once gives the reported
ranges, which describe sensitivity to single blocks. The faithful and exact-differences benchmarks use
the block means of the baseline arm and the same cross-block products. Table~\ref{tab:sdturbo}
gives the values by feature family. Across all features the shared fraction is 64.2\%
(63.9--64.6\% across block deletions), 57.5\% in color and 84.9\% in spatial features but 36.6\%
in texture (36.2--36.9\%); within scenes it is 72.4\% overall and 48.5\% for texture. The
exact-differences value in texture is 46.5\%. SD-Turbo's aligned amplitude is 0.618, its
scene-averaged error 0.917 and its scene-wise error 1.637, and its Monet--Sisley alignment is
0.471. The collection reuses the same references and was analyzed before in other ways.

% Generated by paper/tmlr/build_assets.py; do not edit by hand.
% input reports/painter_cross_cohort_v1/analysis.json sha256 9488dfafa5779b8b065cc756804e67ad64c2841a5c62c2f9b0105f8fef13973e
% input reports/painter_tmlr_diagnostics_v2/analysis.json sha256 bed891f767a98f4b281b6783b1021956f4d2c3ff80becbf6b2693ee45987ef1a
\begin{table}[t]
\caption{Shared fraction (\%) in the SD-Turbo collection by feature family: pooled over scenes, its range over the 25 single-block deletions, computed within scenes, and the faithful and exact-differences benchmarks.}
\label{tab:sdturbo}
\vspace{2pt}
\centering\small
\begin{tabular}{@{}lrcrrr@{}}
\toprule
Features & shared (\%) & block deletions & within scene & faithful & exact diff.\\
\midrule
All 31 & 64.2 & 63.9--64.6 & 72.4 & 92.5 & 67.4\\
Color (11) & 57.5 & 57.3--57.8 & 66.6 & 79.7 & 56.4\\
Spatial (8) & 84.9 & 84.4--85.3 & 81.2 & 96.8 & 83.3\\
Texture (12) & 36.6 & 36.2--36.9 & 48.5 & 89.8 & 46.5\\
\bottomrule
\end{tabular}
\end{table}


\FloatBarrier
\section{The second collection}\label{app:second}

\paragraph{Protocol and amendments.} After the four-painter results were known, a protocol for a
second collection named the two further groups, fixed the tests H1 and H2
(Section~\ref{sec:inference}) and the per-group analysis, and was frozen with the request frame
before any image was generated. The freeze record binds the hashes of the protocol file and of the
analysis code as approved, so the protocol's header still reads ``draft''; the correction of H2 for
panel size is in that frozen code. The plan for the embedding readouts (shared-fraction intervals,
recognition and proximity) was fixed while the images were being collected and before any was
measured. The reference rules were amended twice before the first request. First, the protocol's
minimum of 60 works per painter was withdrawn so that the new collections are treated exactly like
the four-painter collection, which had none; three Hudson River painters and Kirchner fall below it
(Table~\ref{tab:v3-panels}). Second, Wikimedia Commons now serves only standard thumbnail widths,
so the rendered files are 1,280, 1,920 or 3,840 pixels wide (or the original), all reduced to a
512-pixel short side for measurement. The title lexicon keeps every term of the four-painter
lexicon and adds place and landscape terms for the new painters' titles (for example Yosemite,
Catskill, piazza and wheatfield) and exclusion terms for figure and interior titles that those
words would otherwise admit; the terms were fixed from the titles before any image was requested.
A first proposal for the related group, four Russian landscape painters, was replaced before any
reference was acquired, because Wikidata records too few documented oil paintings for three of them.

% Generated by paper/tmlr/build_assets.py; do not edit by hand.
% input data/manifests/painter_specificity_v3/refs-20261002/determination_receipt.json sha256 39e289f1ba44c7d5fe8ae18d59d076a07eb1cefd6d105859b3565bb6576fa982
% input data/manifests/painter_specificity_v3/refs-20261002/measurement_receipt.json sha256 debf37377410a96cd64f1714a277ade667d959714c0c13bbd7af986de0795c16
\begin{table}[t]
\caption{Reference panels of the two further painter groups: Wikidata paintings with a Commons image, and how many remain after each gate, applied in order as for the four-painter panel (single creator, oil on canvas, a collection, an open licence, a short side of at least 1,024 pixels, an outdoor title); then works measured after merging duplicates and excluding unreadable files. Classes are water/built/route/land by title.}
\label{tab:v3-panels}
\vspace{2pt}
\centering\footnotesize
\begin{tabular}{lrrrrrrrrr}
\toprule
Painter & Found & Creator & Medium & Coll. & Rights & Size & Title & Measured & W/B/R/L \\
\midrule
\multicolumn{10}{l}{\emph{Century group}} \\
van Ruisdael & 485 & 470 & 261 & 219 & 218 & 127 & 123 & 123 & 58/31/3/31 \\
Canaletto & 408 & 404 & 342 & 306 & 306 & 166 & 158 & 157 & 120/37/0/1 \\
van Gogh & 887 & 887 & 818 & 638 & 638 & 587 & 255 & 255 & 25/95/12/123 \\
Kirchner & 273 & 273 & 207 & 190 & 187 & 153 & 43 & 41 & 10/17/2/14 \\
\multicolumn{10}{l}{\emph{Hudson River School}} \\
Bierstadt & 481 & 481 & 228 & 138 & 138 & 103 & 92 & 92 & 36/5/0/51 \\
Church & 153 & 152 & 101 & 81 & 81 & 51 & 47 & 47 & 11/3/0/33 \\
Cole & 172 & 172 & 108 & 99 & 99 & 63 & 36 & 36 & 7/7/0/22 \\
Durand & 204 & 204 & 183 & 163 & 163 & 57 & 38 & 37 & 9/2/1/26 \\
\bottomrule
\end{tabular}
\end{table}


\paragraph{References and predictions.} Of 790 works admitted after merging two duplicates, one
file in an unsupported format could not be read and one in an unprofiled non-RGB color space could
not be measured, as for the four-painter collection, leaving 788 (Table~\ref{tab:v3-panels}). The
predictions recorded before collection (Table~\ref{tab:v3-predictions}) use the first collection's
generic outputs; the four-painter rows reproduce this paper's values. The CLIP and CSD checkpoints,
deleted after the first collection's extraction, were downloaded again at the same revisions and
matched their recorded hashes.

% Generated by paper/tmlr/build_assets.py; do not edit by hand.
% input data/manifests/painter_specificity_v3/refs-20261002/predictions.json sha256 65a12c297f2f2e0c2d6d7afbb313e590bf7360001cbe82c5edc5fed20040414f
\begin{table}[t]
\caption{Predictions recorded before the second collection: the reference spread $H$ of each group and the faithful benchmark $N^*/(N^*+H)$ over the six configurations, with September's generic outputs. The four-painter rows reproduce the paper's values.}
\label{tab:v3-predictions}
\vspace{2pt}
\centering\small
\begin{tabular}{lrrrrrr}
\toprule
& \multicolumn{2}{c}{31 features} & \multicolumn{2}{c}{CLIP} & \multicolumn{2}{c}{CSD} \\
\cmidrule(lr){2-3}\cmidrule(lr){4-5}\cmidrule(lr){6-7}
Group & $H$ & faithful (\%) & $H$ & faithful (\%) & $H$ & faithful (\%) \\
\midrule
Four Impressionists & 5.92 & 84.8--95.2 & 0.139 & 86.1--91.2 & 0.325 & 82.4--89.7 \\
Century group & 34.61 & 47.0--73.3 & 0.460 & 67.4--74.8 & 1.048 & 63.4--70.3 \\
Hudson River School & 6.24 & 90.1--96.9 & 0.082 & 92.7--94.3 & 0.142 & 92.2--95.5 \\
\bottomrule
\end{tabular}
\end{table}


\paragraph{Collection.} The six configurations kept the first collection's routes and request
fields; tests in the code check that the no-clause and generic requests are identical to the first
collection's for every configuration and scene, and the gateway listed every route at unchanged
prices on the day of collection. The 1,680 requests were randomized within scene blocks of 120
with a new seed. Two requests were refused by the provider's safety system (HTTP 400, no charge)
for prompts that name Durand and Bierstadt with ordinary landscape scenes; under the protocol they
are recorded and not repeated. Two gateway timeouts (HTTP 503) were repeated successfully. The new
charges were \$72.93. The collector reserves \$5 for every request until its charge is known and
keeps that reservation when a failure reports no charge, so the two timeouts held \$10 that the
bookkeeping never released; the spending ceiling was therefore raised from \$200 to \$220 in a
recorded amendment, which changed no other rule.

% Generated by paper/tmlr/build_assets.py; do not edit by hand.
% input reports/painter_specificity_v3/analysis.json sha256 75ef2a23c04a5c570fd224a0ada396271ad7026ef0657a1908023daca709cbc4
\begin{table}[t]
\caption{The two prespecified tests in every representation (the 31 features in the full view are primary). H1: mean over configurations of the century group's shared fraction minus the Hudson River School's (percentage points), with a 95\% interval from 5,000 paired scene resamples. H2: Spearman correlation over the 28 painter pairs between the name distance and the reference distance, averaged over configurations, with the exact one-sided $p$-value over all 40,320 relabellings of the eight painters (a Mantel test, \citealp{mantel1967detection}); the reference distances are corrected for panel size, and the uncorrected test is shown beside it. In the central square window H1 is unchanged by construction, because the generated images are square and the shared fraction does not use the references. Last columns: reference spread corrected for panel size.}
\label{tab:v3-tests}
\vspace{2pt}
\centering\small\setlength{\tabcolsep}{4pt}
\begin{tabular}{@{}lrcrrrrrr@{}}
\toprule
 & \multicolumn{2}{c}{H1: century $-$ Hudson} & \multicolumn{2}{c}{H2: Mantel} & \multicolumn{2}{c}{H2, uncorrected} & \multicolumn{2}{c}{$H$, noise-corrected}\\
\cmidrule(lr){2-3}\cmidrule(lr){4-5}\cmidrule(lr){6-7}\cmidrule(l){8-9}
Representation & diff. & 95\% & $\rho$ & $p$ & $\rho$ & $p$ & century & Hudson\\
\midrule
31 features & $-$72.7 & [$-$76.5, $-$62.6] & 0.85 & $<$0.001 & 0.84 & $<$0.001 & 33.80 & 4.21\\
31 features, central square & $-$72.7 & [$-$76.5, $-$62.6] & 0.85 & 0.001 & 0.85 & 0.001 & 35.30 & 5.07\\
CLIP & $-$40.1 & [$-$41.4, $-$37.3] & 0.86 & $<$0.001 & 0.86 & $<$0.001 & 0.452 & 0.070\\
CSD & $-$47.4 & [$-$48.9, $-$44.9] & 0.92 & $<$0.001 & 0.92 & $<$0.001 & 1.034 & 0.116\\
\bottomrule
\end{tabular}
\end{table}


\paragraph{Tests and sensitivity.} H1 and H2 hold in every representation, and H2 also holds in the
central square window and without the correction for panel size (Table~\ref{tab:v3-tests}); H1 is
unchanged in the central square window by construction, because the generated images are square
and the shared fraction does not use the references. The per-configuration differences in shared
fraction are large in every configuration. For FLUX.2 Max the adjusted interval is wide because
both of its fractions are unstable under resampling: the prespecified H1 code keeps resamples whose
denominator is not positive, and for FLUX.2 Max with the Hudson River School $N+B\le0$ in 90 of the
5,000 draws, while a fraction lies outside $[0,1]$ in 150 draws for the century group and 237 for
the Hudson River School (Table~\ref{tab:v3-shared}). No other configuration has a non-positive
denominator.

\paragraph{Breakdowns defined after review.} Table~\ref{tab:v3-closeness} splits the two tests
descriptively, under a plan written after a first review of these results; the reviewers had
already computed the point values, and the intervals are new. H2 holds within the century group and
across the groups, and is weak within the Hudson River School, where the reference differences are
small and only six pairs enter. A faithful imitator would show a century-minus-Hudson difference of
30.9 points in the features, so the remaining 41.8 points come from how strongly the generators
separate the names. Table~\ref{tab:v3-readouts} adds the readouts of the embedding plan that the
main text does not quote.

% Generated by paper/tmlr/build_assets.py; do not edit by hand.
% input reports/painter_specificity_v3/analysis.json sha256 75ef2a23c04a5c570fd224a0ada396271ad7026ef0657a1908023daca709cbc4
\begin{table}[t]
\caption{Agreement of the between-name differences with each group's reference differences (31 features, second collection), as in Table~\ref{tab:agreement}: aligned amplitude $\beta$ and error $D$ with unadjusted paired-scene Student intervals, relative size $Q$, alignment ratio, held-out rescaled error, and the share of the centroid proximity gain (Eq.~\ref{eq:prototype-gain}, here in the features) carried by the shared term.}
\label{tab:v3-detail}
\vspace{2pt}
\centering\small\setlength{\tabcolsep}{4pt}
\begin{tabular}{@{}lrcrrrcrr@{}}
\toprule
Configuration & $\beta$ & 95\% & $Q$ & $\beta/\sqrt{Q}$ & $D$ & 95\% & $D_{\mathrm{held}}$ & shared gain (\%)\\
\midrule
\multicolumn{9}{l}{\emph{Century group}}\\
GPT Image 1 & 1.28 & [1.11, 1.45] & 3.13 & 0.72 & 1.58 & [0.74, 2.41] & 0.50 & 100.7\\
GPT Image 2 & 1.14 & [1.06, 1.22] & 2.22 & 0.76 & 0.94 & [0.75, 1.13] & 0.42 & 73.0\\
Flare & 1.29 & [1.21, 1.38] & 2.58 & 0.80 & 0.99 & [0.76, 1.23] & 0.35 & 76.2\\
Sunburst & 1.24 & [1.19, 1.30] & 2.61 & 0.77 & 1.12 & [0.84, 1.40] & 0.41 & 89.1\\
Nano Banana 2 & 1.04 & [0.93, 1.15] & 2.21 & 0.70 & 1.13 & [0.50, 1.77] & 0.53 & 86.2\\
FLUX.2 Max & 0.84 & [0.75, 0.93] & 1.41 & 0.71 & 0.73 & [0.60, 0.85] & 0.50 & 31.5\\
\multicolumn{9}{l}{\emph{Hudson River School}}\\
GPT Image 1 & $-$0.01 & [$-$0.12, 0.09] & 0.22 & $-$0.03 & 1.25 & [0.86, 1.64] & 1.00 & --\\
GPT Image 2 & 0.04 & [$-$0.06, 0.15] & 0.12 & 0.13 & 1.03 & [0.75, 1.31] & 1.05 & 100.2\\
Flare & $-$0.01 & [$-$0.09, 0.06] & 0.61 & $-$0.01 & 1.63 & [1.32, 1.95] & 1.00 & 103.3\\
Sunburst & 0.11 & [0.01, 0.22] & 0.48 & 0.16 & 1.26 & [1.04, 1.47] & 0.98 & 101.4\\
Nano Banana 2 & 0.20 & [$-$0.01, 0.40] & 0.65 & 0.24 & 1.26 & [0.27, 2.25] & 1.14 & 99.7\\
FLUX.2 Max & 0.12 & [$-$0.01, 0.25] & 0.22 & 0.26 & 0.98 & [0.69, 1.26] & 1.14 & 86.7\\
\bottomrule
\end{tabular}
\end{table}

% Generated by paper/tmlr/build_assets.py; do not edit by hand.
% input reports/painter_specificity_v3/analysis.json sha256 75ef2a23c04a5c570fd224a0ada396271ad7026ef0657a1908023daca709cbc4
% input reports/painter_tmlr_diagnostics_v6/analysis.json sha256 16b91e3b0151ab29ef142cdea4a5951a2df40e87d4058ff255eeb20094f4f4af
\begin{table}[t]
\caption{The two further groups in the embeddings: shared fraction with a 95\% interval from 5,000 scene resamples and the faithful benchmark, aligned amplitude, alignment ratio, error, the share of the proximity gain carried by the shared term (Eq.~\ref{eq:prototype-gain}) and four-way recognition (macro accuracy, chance 25\%).}
\label{tab:v3-embeddings}
\vspace{2pt}
\centering\footnotesize\setlength{\tabcolsep}{3.5pt}
\begin{tabular}{@{}lrcrrrrrr@{}}
\toprule
Configuration & shared (\%) & 95\% & faithful (\%) & $\beta$ & $\beta/\sqrt{Q}$ & $D$ & shared gain (\%) & recognized (\%)\\
\midrule
\multicolumn{9}{l}{\emph{CLIP, Century group}}\\
GPT Image 1 & 47.0 & [44.7, 50.4] & 64.6 & 0.58 & 0.57 & 0.86 & 29.5 & 71.9\\
GPT Image 2 & 47.9 & [46.5, 51.9] & 68.6 & 0.69 & 0.61 & 0.89 & 45.5 & 85.4\\
Flare & 51.8 & [50.0, 55.3] & 71.1 & 0.57 & 0.53 & 1.04 & 48.7 & 83.3\\
Sunburst & 52.0 & [50.0, 55.3] & 71.6 & 0.67 & 0.58 & 0.98 & 46.2 & 85.4\\
Nano Banana 2 & 60.7 & [57.3, 64.6] & 74.3 & 0.62 & 0.59 & 0.86 & 49.0 & 81.2\\
FLUX.2 Max & 53.9 & [51.2, 57.6] & 70.9 & 0.63 & 0.59 & 0.89 & 46.2 & 76.0\\
\multicolumn{9}{l}{\emph{CLIP, Hudson River School}}\\
GPT Image 1 & 97.5 & [95.1, 99.6] & 92.1 & 0.14 & 0.25 & 1.02 & 96.2 & 35.4\\
GPT Image 2 & 98.1 & [96.7, 99.0] & 92.3 & 0.10 & 0.21 & 1.01 & 98.2 & 29.2\\
Flare & 89.7 & [85.3, 93.4] & 93.2 & 0.25 & 0.26 & 1.44 & 95.0 & 38.5\\
Sunburst & 88.6 & [83.5, 92.7] & 93.6 & 0.31 & 0.29 & 1.51 & 94.5 & 43.8\\
Nano Banana 2 & 87.4 & [84.0, 90.9] & 94.1 & 0.39 & 0.39 & 1.24 & 92.5 & 50.0\\
FLUX.2 Max & 93.0 & [87.0, 96.3] & 92.8 & 0.18 & 0.24 & 1.19 & 95.7 & 35.4\\
\multicolumn{9}{l}{\emph{CSD, Century group}}\\
GPT Image 1 & 46.2 & [44.8, 48.2] & 66.2 & 0.75 & 0.67 & 0.74 & 6.5 & 76.0\\
GPT Image 2 & 40.6 & [39.8, 42.7] & 62.8 & 0.95 & 0.76 & 0.67 & 15.4 & 89.6\\
Flare & 44.3 & [43.2, 46.5] & 69.2 & 0.83 & 0.69 & 0.77 & 23.4 & 83.3\\
Sunburst & 44.2 & [43.0, 46.6] & 68.6 & 0.93 & 0.74 & 0.72 & 25.9 & 85.4\\
Nano Banana 2 & 55.4 & [52.1, 58.6] & 70.6 & 0.74 & 0.68 & 0.71 & 40.7 & 76.0\\
FLUX.2 Max & 51.7 & [48.9, 54.9] & 63.6 & 0.75 & 0.71 & 0.61 & 32.6 & 72.9\\
\multicolumn{9}{l}{\emph{CSD, Hudson River School}}\\
GPT Image 1 & 97.3 & [96.1, 98.2] & 93.2 & 0.12 & 0.21 & 1.08 & 95.1 & 31.2\\
GPT Image 2 & 98.3 & [97.2, 99.0] & 93.6 & 0.11 & 0.23 & 1.02 & 98.3 & 32.3\\
Flare & 93.2 & [91.8, 94.4] & 95.3 & 0.19 & 0.21 & 1.45 & 96.5 & 34.4\\
Sunburst & 92.3 & [90.4, 93.9] & 95.4 & 0.30 & 0.31 & 1.37 & 95.4 & 45.8\\
Nano Banana 2 & 92.3 & [89.4, 94.6] & 95.4 & 0.38 & 0.40 & 1.15 & 94.3 & 44.8\\
FLUX.2 Max & 93.2 & [88.5, 96.7] & 92.1 & 0.20 & 0.35 & 0.94 & 94.9 & 38.5\\
\bottomrule
\end{tabular}
\end{table}

% Generated by paper/tmlr/build_assets.py; do not edit by hand.
% input reports/painter_tmlr_diagnostics_v7/analysis.json sha256 385f089c29f6853d79c9560dd80243b6ac331583334f03801e4e11b6722aac74
\begin{table}[t]
\caption{Two breakdowns of the prespecified tests, defined after review (descriptive). Left: the H2 Spearman correlation, averaged over configurations, over all 28 pairs, the 6 pairs within each group and the 16 pairs across the groups. Right: H1's mean century-minus-Hudson difference, observed and for a faithful imitator (closeness alone), and the excess, the observed minus the faithful difference, with 95\% intervals from 5,000 paired scene resamples.}
\label{tab:v3-closeness}
\vspace{2pt}
\centering\footnotesize\setlength{\tabcolsep}{3.5pt}
\begin{tabular}{@{}lrrrrrrcrc@{}}
\toprule
 & \multicolumn{4}{c}{H2 by pair type (Spearman)} & \multicolumn{5}{c}{Century $-$ Hudson shared fraction (points)}\\
\cmidrule(lr){2-5}\cmidrule(l){6-10}
Representation & all & century & Hudson & across & obs. & faithful & 95\% & excess & 95\%\\
\midrule
31 features & 0.85 & 0.78 & 0.09 & 0.74 & $-$72.7 & $-$30.9 & [$-$32.9, $-$26.6] & $-$41.8 & [$-$47.2, $-$34.9]\\
CLIP & 0.86 & 0.73 & 0.49 & 0.76 & $-$40.1 & $-$22.9 & [$-$23.5, $-$20.4] & $-$17.3 & [$-$19.2, $-$15.7]\\
CSD & 0.92 & 0.88 & 0.27 & 0.85 & $-$47.4 & $-$27.3 & [$-$28.5, $-$25.0] & $-$20.0 & [$-$21.7, $-$18.7]\\
\bottomrule
\end{tabular}
\end{table}

% Generated by paper/tmlr/build_assets.py; do not edit by hand.
% input reports/painter_specificity_v3/analysis.json sha256 75ef2a23c04a5c570fd224a0ada396271ad7026ef0657a1908023daca709cbc4
% input reports/painter_tmlr_diagnostics_v6/analysis.json sha256 16b91e3b0151ab29ef142cdea4a5951a2df40e87d4058ff255eeb20094f4f4af
\begin{table}[t]
\caption{Further embedding readouts of the second collection, fixed before its images were measured. Top: recognition among all eight painters (macro accuracy, \%, chance 12.5\%). Bottom, across the six configurations: the Spearman correlation between the alignment ratio and four-way recognition, and the Pearson correlations of the proximity gain with its two terms (Eq.~\ref{eq:prototype-gain}).}
\label{tab:v3-readouts}
\vspace{2pt}
\centering\small
\begin{tabular}{@{}lrrrr@{}}
\toprule
 & \multicolumn{2}{c}{CLIP} & \multicolumn{2}{c}{CSD}\\
\cmidrule(lr){2-3}\cmidrule(l){4-5}
 & century & Hudson & century & Hudson\\
\midrule
GPT Image 1 & 65.6 & 34.4 & 64.6 & 31.2\\
GPT Image 2 & 78.1 & 29.2 & 87.5 & 32.3\\
Flare & 62.5 & 38.5 & 76.0 & 34.4\\
Sunburst & 69.8 & 43.8 & 83.3 & 45.8\\
Nano Banana 2 & 74.0 & 50.0 & 75.0 & 44.8\\
FLUX.2 Max & 68.8 & 35.4 & 69.8 & 38.5\\
\midrule
Spearman, alignment ratio and four-way recognition & 0.29 & 0.94 & 0.64 & 0.66\\
Correlation of the gain with its shared term & 0.97 & 0.99 & 0.83 & 1.00\\
Correlation of the gain with its painter-specific term & 0.75 & 0.54 & 0.34 & 0.54\\
\bottomrule
\end{tabular}
\end{table}


\paragraph{Change between the collections.} The no-clause and generic requests of the two
collections are identical. Relative to the first collection's repeat noise, the squared distance
between their scene means, estimated without repeat noise, ranges from $-$0.61 to 0.48 over
configurations and arms in the 31 features and from $-$0.29 to 0.24 in the embeddings
(Table~\ref{tab:v3-drift}); values near zero mean no change beyond the noise between two repeats.

% Generated by paper/tmlr/build_assets.py; do not edit by hand.
% input reports/painter_specificity_v3/analysis.json sha256 75ef2a23c04a5c570fd224a0ada396271ad7026ef0657a1908023daca709cbc4
\begin{table}[t]
\caption{Change between the two collections for identical requests: the squared distance between the September and October scene means of the no-clause and generic arms, estimated without repeat noise, divided by September's repeat noise (0 means no change; values near or above 1 mean a change as large as two repeats differ).}
\label{tab:v3-drift}
\vspace{2pt}
\centering\small
\begin{tabular}{@{}lrrrrrr@{}}
\toprule
 & \multicolumn{2}{c}{31 features} & \multicolumn{2}{c}{CLIP} & \multicolumn{2}{c}{CSD}\\
\cmidrule(lr){2-3}\cmidrule(lr){4-5}\cmidrule(l){6-7}
Configuration & none & generic & none & generic & none & generic\\
\midrule
GPT Image 1 & $-$0.15 & 0.22 & $-$0.25 & 0.04 & $-$0.29 & 0.18\\
GPT Image 2 & 0.17 & 0.24 & 0.22 & 0.03 & 0.22 & $-$0.00\\
Flare & $-$0.15 & 0.15 & 0.03 & 0.07 & $-$0.15 & 0.01\\
Sunburst & 0.48 & $-$0.13 & $-$0.05 & $-$0.09 & 0.03 & $-$0.03\\
Nano Banana 2 & $-$0.27 & $-$0.14 & $-$0.06 & $-$0.18 & $-$0.14 & $-$0.14\\
FLUX.2 Max & $-$0.10 & $-$0.61 & 0.14 & 0.24 & $-$0.09 & 0.16\\
\bottomrule
\end{tabular}
\end{table}


\FloatBarrier
\section{Evidence and replay}\label{app:replay}

Every analysis in this paper reads retained feature vectors or embeddings and writes a versioned
result file whose inputs are recorded by SHA-256. All tables and Figures~\ref{fig:benchmark}, \ref{fig:v3-pairs}, \ref{fig:pairs}
and~\ref{fig:projection} are generated from those result files, and Figure~\ref{fig:schematic} by
the same script; Figure~\ref{fig:examples} is a static image panel whose bytes and selection manifest are
checked by hash. The check script verifies the input hashes, regenerates every generated table and
figure byte for byte, and confirms that each registered number quoted in the text equals its
recomputed value and appears in the sentence recorded for it. These checks establish that the
reported numbers follow from the retained measurements; they do not re-extract features from
pixels.
